\documentclass[11pt,a4paper]{amsart}
\usepackage[T1]{fontenc}
\usepackage{lmodern}
\usepackage{amssymb}
\usepackage{microtype}
\usepackage[hyphens]{url}
\usepackage[colorlinks=true,linkcolor=black,citecolor=black,urlcolor=black]{hyperref}
\usepackage[margin=2.5cm]{geometry}

\hypersetup{
  pdftitle={Three colours for the plane over a number field},
  pdfauthor={Sergi Gal\'an},
  pdfsubject={The Hadwiger-Nelson problem over number fields},
  pdfkeywords={Hadwiger-Nelson problem, chromatic number of the plane, unit-distance graph, number field,
  local-global principle, circular chromatic number}}

\newtheorem{theorem}{Theorem}
\renewcommand{\thetheorem}{\Alph{theorem}}
\newtheorem*{thmW}{Theorem W$^+$}
\newtheorem{lemma}{Lemma}
\newtheorem{corollary}{Corollary}
\newtheorem{proposition}{Proposition}
\theoremstyle{remark}
\newtheorem*{remark}{Remark}

\newcommand{\Q}{\mathbb{Q}}
\newcommand{\R}{\mathbb{R}}
\newcommand{\Z}{\mathbb{Z}}
\newcommand{\C}{\mathbb{C}}
\newcommand{\F}{\mathbb{F}}
\newcommand{\Cay}{\operatorname{Cay}}
\newcommand{\Tr}{\operatorname{Tr}}
\newcommand{\cA}{\mathcal{A}}

\title{Three colours for the plane over a number field}
\author{Sergi Gal\'an}
\thanks{These results were obtained with the assistance of an AI system (Claude, by Anthropic), which found the
proofs, wrote them and the programs that check them, and helped to draft this note. Separate AI agents refereed the
proofs, with their own programs. No mathematician has checked the note yet. The author is not a professional
mathematician and submits it for independent verification.}
\date{Draft, 4 October 2026}
\subjclass[2020]{Primary 05C15, 52C10; Secondary 11R04, 11S15, 05C25}
\keywords{Hadwiger--Nelson problem, chromatic number of the plane, unit-distance graph, number field, local--global
principle, circular chromatic number}

\begin{document}

\begin{abstract}
For a field $F$ of characteristic $0$ let $\chi(F^2)$ be the chromatic number of the graph on $F^2$ that joins $x$
and $y$ when $(x_1-y_1)^2+(x_2-y_2)^2=1$. We prove that for every number field $F$, $\chi(F^2)\le3$ if and only if
some prime of $F$ above $2$ ramifies in $F(i)$ or some prime of $F$ above $3$ has residue degree $1$. The ``if'' part
is classical (colourings through a residue field); the ``only if'' part extends to all number fields our earlier
result that $\chi(\Q(\sqrt d)^2)\ge4$ for $d\equiv11\pmod{12}$. It answers, for three colours, the question whether
the chromatic number of the plane over a number field is decided by a single completion, and it classifies the
finite extensions $K$ of $\Q_p$: $\chi(K^2)$ is $2$, $3$ or at least $4$ according to explicit local conditions. We
also give an elementary family (every real field containing $\sqrt a$ and $\sqrt b$ with $a\equiv2\pmod 3$ and
$b\equiv7\pmod8$ needs four colours), new examples such as $\chi(\Q(\sqrt2,\sqrt7)^2)=4$, and, as far as we know,
the first non-integral value of the circular chromatic number of the plane over a number field:
$\chi_c(\Q(\sqrt{11})^2)=\chi_c(\Q(\sqrt{35})^2)=7/2$ (and $\chi_c(\Q_7^2)=7/2$), although these planes have chromatic
number $4$. Locally constant colourings of one completion give no upper bound on $\chi_c$ below $4$ other than $2$, $3$ and
$7/2$. Finally, $\chi_c(F^2)\ge56/17$ whenever $\chi(F^2)\ge4$, so the circular chromatic number of the plane over a
number field never lies in $(3,56/17)$; with a computer-assisted step, not even in $(3,3.3315]$.
\end{abstract}

\maketitle

\section{Introduction}

Let $F$ be a field of characteristic $0$. The \emph{plane over $F$} is the graph on $F^2$ in which $x$ and $y$ are
adjacent when $(x_1-y_1)^2+(x_2-y_2)^2=1$, and $\chi(F^2)$ is its chromatic number. For $F=\R$ this is the
Hadwiger--Nelson problem, and $5\le\chi(\R^2)\le7$ \cite{deGrey}. Over number fields the question goes back to
Woodall ($\chi(\Q^2)=2$ \cite{Woodall}) and was studied by Johnson \cite{Johnson1987}, Fischer \cite{Fischer1990},
Benda and Perles \cite{BP}, Moorhouse \cite{Moorhouse} and Madore \cite{Madore}. All upper bounds below $7$ known to us come from a colouring
that is locally constant at one completion $F_v$: reduction modulo a power of a prime. Writing $L=F(i)$, the unit
vectors of $F^2$ are the elements of norm one of $L$, and at a place where $x^2+y^2$ is anisotropic they reduce into
a finite group of roots of unity of the residue field. Two such reductions give colourings with at most three colours:
\begin{itemize}
\item[(a)] at a prime of $F$ above $2$ that ramifies in $F(i)$, every unit vector is $\equiv1$ modulo the maximal
ideal, which gives a $2$-colouring (Johnson; Fischer \cite[Thm.~8]{Fischer1990}; Moorhouse \cite{Moorhouse};
\cite[Theorem~A$'$]{MM});
\item[(b)] at a prime of $F$ above $3$ of residue degree $1$, the unit vectors reduce into $\mu_4\subset\F_9$, and
$a+bi\mapsto a+b$ gives a $3$-colouring (Madore \cite[Cor.~3.4]{Madore}; for real quadratic fields Moorhouse
\cite[Lemma~8.2 and Cor.~8.3]{Moorhouse} and Fischer \cite[Thm.~9]{Fischer1990}).
\end{itemize}
In \cite{TC} we proved that (a) is also necessary for $\chi(F^2)=2$. Here we prove the same at three colours.

\setcounter{theorem}{1}
\begin{theorem}\label{thm:B}
Let $F$ be a number field. Then $\chi(F^2)\le3$ if and only if (a) or (b) holds; and $\chi(F^2)\le2$ if and only
if (a) holds.
\end{theorem}

For real quadratic fields $\Q(\sqrt d)$, $d$ squarefree, (a) holds iff $d\not\equiv3\pmod4$ and (b) iff
$d\not\equiv2\pmod3$, so Theorem~\ref{thm:B} contains the main result of \cite{FC} (Theorem~1 and Corollary~2
there): $\chi(\Q(\sqrt d)^2)\ge4$ iff $d\equiv11\pmod{12}$. Whether (a) or (b) holds depends only on $F\otimes\Q_2$ and $F\otimes\Q_3$, and is computed from
a defining polynomial (the relative discriminant of $F(i)/F$, and the residue degrees above $3$). By the theorem of
de Bruijn and Erd\H{o}s \cite{dBE}, $\chi(F^2)\ge4$ if and only if some finite unit-distance graph in $F^2$ needs four
colours, and Theorem~\ref{thm:B} says that this happens unless one of the completions of $F$ has a locally constant
$3$-colouring: a local--global principle for three colours. The same proof decides the local fields
(Corollary~\ref{cor:local}): for a finite extension $K$ of $\Q_p$, $\chi(K^2)=2$ if $p=2$ and $K(i)/K$ is ramified,
$\chi(K^2)=3$ if $p=3$ and $K$ has residue degree $1$, and $\chi(K^2)\ge4$ otherwise; for instance the plane over the
unramified cubic extension of $\Q_3$ needs at least four colours, although $x^2+y^2$ is anisotropic there, as over
$\Q_3$.

Two families of fields were known to need four colours: those containing $\sqrt d$ with $d\equiv11\pmod{12}$
\cite{FC}, and those containing $\sqrt3$ and $\sqrt q$ with $q\equiv2\pmod3$, by chains of unit rhombi
(\cite{Repo}, \nolinkurl{notes/local_colourings.md}, Section~11, Theorem~5). Theorem
\ref{thm:B} adds, for example, $\Q(\sqrt2,\sqrt7)$, $\Q(\sqrt2,\sqrt{31})$ and $\Q(\sqrt2,\sqrt{55})$, which contain
no unit triangle and whose quadratic subfields all have $\chi\le3$, and $\Q(2\cos\frac{2\pi}7,\sqrt7)$, of degree
$6$, where $3$ has residue degree $3$. For many fields an explicit proof is short (Proposition~\ref{prop:two}), and
for $\Q(\sqrt2,\sqrt7)$ we give a certificate checked by computer that does not depend on the general argument.

\begin{theorem}\label{thm:C}
$\chi_c(\Q(\sqrt{11})^2)=\chi_c(\Q(\sqrt{35})^2)=\chi_c(\Q_7^2)=7/2$, while the three chromatic numbers are $4$.
Moreover no plane over a number field has circular chromatic number strictly between $2$ and $3$.
\end{theorem}

\begin{theorem}\label{thm:D}
Let $F$ be a number field with $\chi(F^2)\ge4$. Then $\chi_c(F^2)\ge56/17\approx3.294$. With a computer-assisted
step \textup{(}exact computations, made by three programs written separately\textup{)},
$\chi_c(F^2)\ge10^7/3001611>3.3315$.
\end{theorem}

So the circular chromatic number of the plane over a number field is $2$, $3$ or at least $56/17$ (larger than
$3.3315$, with the computer-assisted step); we do not know whether a value in $(3,7/2)$ occurs
(Section~\ref{sec:questions}).

The tools are those of \cite{FC}: a characterisation of $3$-colourings of Cayley graphs of abelian groups by
characters (Theorem~W$^+$, Section~\ref{sec:prelim}), and the exact shape of the characters along the rational rotations whose
denominators are powers of $5$. These rotations act as a probe: along them, a character that keeps every unit vector
in $[1/3,2/3]$ is, up to a bounded error, either a $2$-adic ``half'' character or a $3$-adic character of level one,
and a compactness argument turns this into the dichotomy between the places above $2$ and those above $3$. For
Theorem~\ref{thm:D} we show that the probe keeps this shape, with a small error also for type~q, when $[1/3,2/3]$ is
replaced by $[\theta,1-\theta]$ with $\theta>17/56$ (Proposition~\ref{prop:probe}).

Sections~\ref{sec:prelim}--\ref{sec:three} prove Theorem~\ref{thm:B}; Section~\ref{sec:cor} draws the consequences
(circular colourings below $3$, fields split above $2$ and $3$, local fields); Section~\ref{sec:two} gives the
explicit family; Section~\ref{sec:circ} proves Theorem~\ref{thm:C} and shows that the restrictions to $F^2$ of locally
constant colourings of a completion $F_v^2$ give no upper bound on $\chi_c(F^2)$ below $4$ other than $2$, $3$ and
$7/2$ (Corollary~\ref{cor:localvalues}); Section~\ref{sec:gap} proves Theorem~\ref{thm:D}; Section~\ref{sec:checks}
lists the computer checks, and Section~\ref{sec:questions} asks
whether $2$, $3$ and $7/2$ are the only values below $4$ for number fields. The
programs and certificates are in \cite{Repo}.

\section{Preliminaries}\label{sec:prelim}

\emph{Characters and colourings.} For an abelian group $\Gamma$ and $S=-S\subseteq\Gamma$, $\Cay(\Gamma,S)$ is the
graph on $\Gamma$ in which $x\sim y$ when $y-x\in S$. A \emph{character} of $\Gamma$ is a homomorphism
$\xi:\Gamma\to\R/\Z$, and $\|t\|$ is the distance from $t$ to the nearest integer. For integers $p\ge2q\ge1$, the
circular clique $K_{p/q}$ has vertex set $\Z/p$, with $a\sim b$ when $q\le(b-a\bmod p)\le p-q$, and the circular
chromatic number $\chi_c(G)$ is the infimum of the $p/q$ such that $G$ has a homomorphism to $K_{p/q}$.

\begin{thmW}[{\cite[Theorem~1]{Winding}}]
Let $p\ge2q\ge2$ with $p<4q$. Then $\Cay(\Gamma,S)$ has a homomorphism to $K_{p/q}$ if and only if some character
$\xi$ of $\Gamma$ satisfies $\xi(s)\in[q/p,1-q/p]+\Z$ for every $s\in S$. In particular $\Cay(\Gamma,S)$ is
$3$-colourable if and only if some character maps $S$ into $[1/3,2/3]+\Z$.
\end{thmW}

The ``if'' direction is classical (colour $x$ by $\lfloor p\,\xi(x)\rfloor$); the ``only if'' direction is the
content of \cite{Winding}, where it is also proved in Lean for every abelian group $\Gamma$ and finite $S$. We apply
it to the groups $\Gamma=\Z U$ spanned by finite sets $U$ of unit vectors, and in Section~\ref{sec:circ} to finite groups
$O_w/\mathfrak m^k$.

\emph{The plane as a Cayley graph.} Let $F$ be a number field with $i\notin F$, and $L=F(i)$. The map
$(x,y)\mapsto x+iy$ identifies $F^2$ with $L$, and $x^2+y^2=z\bar z$ for $z=x+iy$, where $z\mapsto\bar z$ is the
non-trivial automorphism of $L/F$. So the unit vectors form the group $T=\{z\in L:z\bar z=1\}$, and the plane over
$F$ is $\Cay(L,T)$. As a $\Q(i)$-vector space $L$ has dimension $n=[F:\Q]$, and $\Tr=\Tr_{L/\Q(i)}$ gives a
non-degenerate pairing $L\times L\to\Q(i)$.

\emph{The probe.} Let $N=5^k$, $k\ge1$, let $G_N$ be the set of $\gamma\in\Q(i)$ with $\gamma\bar\gamma=1$ and
$N\gamma\in\Z[i]$ (the rational rotations whose denominator divides $N$), $I=[1/3,2/3]$, and
\[
S_N=\{c\in\C:\ \operatorname{Re}(\bar c\gamma)\in I+\Z\text{ for every }\gamma\in G_N\}.
\]
Let $E_c=\frac{1+i}2+\Z[i]$, $E_q=\{\frac{\alpha+\beta i}3:\alpha,\beta\in\{1,2\}\}+\Z[i]$, $E=E_c\cup E_q$, and
$r=\sqrt{10}/18$. Elements of $E_c$ have \emph{type~c}, those of $E_q$ \emph{type~q}.

\begin{proposition}[{\cite[Proposition~9]{FC}}]\label{prop:shape}
Every $c\in S_N$ can be written $c=N\varepsilon+x$, either with $\varepsilon\in E_c$ and $|x|\le r$, or with
$\varepsilon\in E_q$ and $x=0$.
\end{proposition}

\begin{lemma}[exact relations; cf.\ {\cite[Lemma~10]{FC}}]\label{lem:exact}
Let $\sigma_j=N\varepsilon_j+x_j\in S_N$ be as in Proposition~\ref{prop:shape} \textup{(}$1\le j\le\ell$\textup{)},
and let $\mu_1,\dots,\mu_\ell\in\Z[i]$ satisfy $\sum_j\mu_j\sigma_j=0$ and $N>6r\sum_j|\mu_j|$. Then
$\sum_j\mu_j\varepsilon_j=0$.
\end{lemma}

\begin{proof}
$\sum_j\mu_jN\varepsilon_j=-\sum_j\mu_jx_j$. The left side lies in $\frac N6\Z[i]$, because $6E\subseteq\Z[i]$,
and the nonzero elements of $\frac N6\Z[i]$ have modulus at least $\frac N6$; the right side has modulus at most
$r\sum_j|\mu_j|<\frac N6$. So both sides vanish.
\end{proof}

Only two properties of $E$ are used below. Let $\F_9=\Z[i]/3$ and $A'=\{\pm1\pm i\}\subset\F_9$, the set of
non-squares of $\F_9$ (equivalently, of the $y\in\F_9$ with $y^4=-1$). Then
\begin{equation}\label{eq:types}
\varepsilon\in E_c\ \Rightarrow\ v_{1+i}(\varepsilon)=-1\text{ and }v_3(\varepsilon)\ge0;\qquad
\varepsilon\in E_q\ \Rightarrow\ v_{1+i}(\varepsilon)\ge0\text{ and }3\varepsilon\bmod3\in A'.
\end{equation}
In particular every $\varepsilon\in E$ has $v_{1+i}(\varepsilon)\ge-1$ and $v_3(\varepsilon)\ge-1$, with equality
in the second exactly for type~q.

\section{From colourings to a linear problem over \texorpdfstring{$\Q(i)$}{Q(i)}}\label{sec:linear}

\begin{lemma}\label{lem:B1}
Let $F$ be a number field with $i\notin F$, and suppose that $F^2$ is $3$-colourable. Then for every finite
$V\subset T$ there is $\beta\in L$ with $\Tr(\beta v)\in E$ for every $v\in V$.
\end{lemma}

\begin{proof}
We may enlarge $V$, so we assume that it contains a $\Q(i)$-basis $b_1,\dots,b_n$ of $L$. Such a basis exists in
$T$: the $\Q(i)$-span $W$ of $T$ is a subring of $L$, as $T$ is a group; being finite-dimensional over $\Q(i)$ and
contained in a field, it is a field. For $t\in F$ the element $z=(t+i)/(t-i)$ lies in $T$, so $W$ contains
$(z-1)^{-1}$ and $t=i(z+1)/(z-1)$; hence $W=L$.

For $v\in V$ write $D_vv=\sum_j\mu_{v,j}b_j$ with $D_v\in\Z\setminus\{0\}$ and $\mu_{v,j}\in\Z[i]$, and choose
$N=5^k$ with $N>6r\max_v(|D_v|+\sum_j|\mu_{v,j}|)$. The set $U=G_NV$ is finite, symmetric ($-1\in G_N$) and
consists of unit vectors, so $\Cay(\Z U,U)$ is a subgraph of $\Cay(L,T)$ and is $3$-colourable. By
Theorem~W$^+$ there is a character $\xi$ of $\Z U$ with $\xi(U)\subseteq I+\Z$. Extend $\xi$ to $L$ (the group
$\R/\Z$ is divisible) and restrict it to a lattice $\Lambda=\frac1M\bigoplus_j\Z[i]f_j\supseteq U$, where
$f_1,\dots,f_n$ is a $\Q$-basis of $F$, so that $L=\bigoplus_j\Q(i)f_j$. A character of $\Lambda\cong\Z^{2n}$ lifts
to a real-linear form, and every real-linear form on $\C^n$ is the real part of a $\C$-linear one; so there is a
$\Q(i)$-linear map $\varphi:L\to\C$ with $\xi(w)=\operatorname{Re}\varphi(w)\bmod1$ for $w\in\Lambda$. For $v\in V$
and $\gamma\in G_N$ we get $\xi(\gamma v)=\operatorname{Re}(\gamma\varphi(v))$, so $\overline{\varphi(v)}\in S_N$,
and $\varphi(v)\in S_N$ because $G_N$ is closed under conjugation. Write $\varphi(v)=N\varepsilon_v+x_v$ as in
Proposition~\ref{prop:shape}. Applying $\varphi$ to $D_vv-\sum_j\mu_{v,j}b_j=0$ and Lemma~\ref{lem:exact} give
$D_v\varepsilon_v=\sum_j\mu_{v,j}\varepsilon_{b_j}$. Let $\beta\in L$ be the element with
$\Tr(\beta b_j)=\varepsilon_{b_j}$ for all $j$. Then $D_v\Tr(\beta v)=\sum_j\mu_{v,j}\varepsilon_{b_j}=D_v\varepsilon_v$,
so $\Tr(\beta v)=\varepsilon_v\in E$ for every $v\in V$.
\end{proof}

Lemma~\ref{lem:B1} is also a practical test: if for some finite $V\subset T$ no such $\beta$ exists, then $F^2$ is
not $3$-colourable. Whether $\beta$ exists is a finite computation (Section~\ref{sec:checks}).

\section{Localisation at \texorpdfstring{$2$}{2} and \texorpdfstring{$3$}{3}}\label{sec:local}

For a prime $p$ let $L_p=L\otimes_\Q\Q_p=\prod_{w\mid p}L_w$, and for a place $u$ of $F$ let
$L_u=L\otimes_FF_u$ and $T(F_u)=\{z\in L_u:z\bar z=1\}$. Put $T_p=\prod_{u\mid p}T(F_u)$. Since $2$ ramifies and $3$
is inert in $\Q(i)$, $\Q(i)\otimes\Q_2=\Q_2(i)$ and $\Q(i)\otimes\Q_3=\Q_9$, the unramified quadratic extension of
$\Q_3$; the trace extends to $\Tr:L_2\to\Q_2(i)$ and $\Tr:L_3\to\Q_9$. Let $O_2=\Z_2[i]$, $O_3=\Z_9$ (the integers of
$\Q_9$), and $\cA=\{\frac{\pm1\pm i}3\}+O_3$. By \eqref{eq:types}, $E_c\subseteq(1+i)^{-1}O_2^\times\cap O_3$ and
$E_q\subseteq O_2\cap\cA$.

\begin{lemma}[weak approximation for $T$]\label{lem:WA}
For every finite set $\Sigma$ of places of $F$, $T$ is dense in $\prod_{u\in\Sigma}T(F_u)$.
\end{lemma}

\begin{proof}
At a place $u$ with $i\notin F_u$, $t\mapsto(t+i)/(t-i)$ is a homeomorphism from $F_u$ onto $T(F_u)\setminus\{1\}$,
with inverse $z\mapsto i(z+1)/(z-1)$. At a place with $i\in F_u$, $L_u\cong F_u\times F_u$,
$T(F_u)=\{(s,s^{-1})\}$, and $t\mapsto(s,s^{-1})$ with $s=(t+\iota)/(t-\iota)$, $\iota^2=-1$, is a homeomorphism
from $F_u\setminus\{\pm\iota\}$ onto $T(F_u)\setminus\{1\}$. These maps send $F$ into $T$, $F$ is dense in
$\prod_{u\in\Sigma}F_u$, and $T(F_u)$ has no isolated points.
\end{proof}

\begin{lemma}\label{lem:B2}
If $F^2$ is $3$-colourable and $i\notin F$, then one of the following holds:
\begin{itemize}
\item[(c)] there is $\beta_2\in L_2$ with $v_{1+i}(\Tr(\beta_2z))=-1$ for every $z\in T_2$;
\item[(q)] there is $\beta_3\in L_3$ with $\Tr(\beta_3z)\in\cA$ for every $z\in T_3$.
\end{itemize}
\end{lemma}

\begin{proof}
Fix a $\Q(i)$-basis $b_1,\dots,b_n\subset T$ of $L$ (proof of Lemma~\ref{lem:B1}), and let
\[
K_2=\{\beta\in L_2:\Tr(\beta b_j)\in(1+i)^{-1}O_2\ \forall j\},\qquad K_3=\{\beta\in L_3:\Tr(\beta b_j)\in\tfrac13O_3\ \forall j\}.
\]
They are compact, since $\beta\mapsto(\Tr(\beta b_j))_j$ is a linear isomorphism onto $\Q_2(i)^n$, resp.\
$\Q_9^n$. On the compact space $\Omega=\{\mathrm c,\mathrm q\}^T\times K_2\times K_3$ consider, for $v\in T$, the
condition $P_v(\tau,\beta_2,\beta_3)$:
\[
\tau(v)=\mathrm c,\ v_{1+i}(\Tr(\beta_2v))=-1,\ \Tr(\beta_3v)\in O_3;\quad\text{or}\quad
\tau(v)=\mathrm q,\ \Tr(\beta_2v)\in O_2,\ \Tr(\beta_3v)\in\cA.
\]
Each $P_v$ defines a clopen subset of $\Omega$. For a finite $V\supseteq\{b_1,\dots,b_n\}$, Lemma~\ref{lem:B1} gives
$\beta\in L$; its images in $L_2$ and $L_3$ lie in $K_2$ and $K_3$, and with $\tau(v)$ the type of $\Tr(\beta v)$ the
triple satisfies $P_v$ for every $v\in V$, by \eqref{eq:types}. So the closed sets $\{P_v\text{ for all }v\in V\}$
have the finite intersection property, and some $(\tau,\beta_2,\beta_3)\in\Omega$ satisfies $P_v$ for every $v\in T$.

For $z\in T_2$ let $a(z)$ be the statement $v_{1+i}(\Tr(\beta_2z))=-1$ and $a'(z)$ the statement
$\Tr(\beta_2z)\in O_2$; for $y\in T_3$ let $b(y)$ be $\Tr(\beta_3y)\in O_3$ and $b'(y)$ be $\Tr(\beta_3y)\in\cA$.
They are locally constant, $a$ and $a'$ exclude each other, and so do $b$ and $b'$. The set
\[
Z=\{(z,y)\in T_2\times T_3:\ (a(z)\wedge b(y))\vee(a'(z)\wedge b'(y))\}
\]
is clopen and contains the diagonal image of $T$, which is dense by Lemma~\ref{lem:WA}; so $Z=T_2\times T_3$. If
$a(z_0)$ holds for some $z_0$, then $b(y)$ holds for every $y$, so $b'$ holds nowhere, and $a(z)$ holds for every
$z$: this is (c). Otherwise $a'$ holds everywhere, and then $b'$ holds everywhere: this is (q).
\end{proof}

\emph{Case (c).} For $z\in L$ let $g(z)=(1+i)\Tr(\beta_2z)\in\Q_2(i)$. Fix a representative $\rho$ of each class of
$\Q_2(i)$ modulo $O_2$, and colour $z$ by $(g(z)-\rho)\bmod(1+i)\in O_2/(1+i)=\F_2$, where $\rho$ represents the
class of $g(z)$. A unit step $v\in T$ adds $g(v)\in O_2^\times$ to $g(z)$: it keeps the class modulo $O_2$ and
changes the residue modulo $1+i$. This is a proper $2$-colouring of $L\cong F^2$, so by \cite[Theorem~1]{TC} some
prime of $F$ above $2$ ramifies in $F(i)$.

\section{Case (q): the places above \texorpdfstring{$3$}{3}}\label{sec:three}

Let $u$ be a place of $F$ above $3$, $f=f_u$ its residue degree, $q=3^f$, and $g_u(z)=\Tr_{L_u/\Q_9}(\beta_uz)$,
where $\beta_3=(\beta_u)_u$; then $\Tr(\beta_3y)=\sum_ug_u(y_u)$ for $y=(y_u)\in T_3$.

\emph{Even $f$.} Then $\F_9$ lies in the residue field of $F_u$, so $i\in F_u$, $L_u\cong F_u\times F_u$ and
$T(F_u)=\{(s,s^{-1}):s\in F_u^\times\}$. Write $g_u(s,t)=\ell_1(s)+\ell_2(t)$ with $\Q_3$-linear
$\ell_1,\ell_2:F_u\to\Q_9$. In case (q), $g_u(y_u)$ stays in a bounded set when $y_u$ varies and the other
components are fixed. If $\ell_1(s_0)\ne0$, then $g_u(3^{-m}s_0,3^ms_0^{-1})=3^{-m}\ell_1(s_0)+3^m\ell_2(s_0^{-1})$
is unbounded as $m\to\infty$; so $\ell_1=0$, likewise $\ell_2=0$, and $g_u=0$.

\emph{Odd $f$.} Then $L_u=F_u(i)=F_u\Q_9$ is a field, unramified of degree $2$ over $F_u$, with residue field
$\F_{q^2}$; its non-trivial automorphism $\sigma$ acts on $\Q_9$ as the Frobenius ($\sigma(i)=-i$). The group
$T(F_u)$ lies in $O_{L_u}^\times$, contains $\mu_4\subset\Q_9$ and the Teichm\"uller lifts of
$\mu_{q+1}\subset\F_{q^2}^\times$ (as $\zeta\sigma(\zeta)=\zeta^{1+q}$), and reduces onto $\mu_{q+1}$. Since $g_u$
is $\Q_9$-linear, $g_u(\epsilon z)=\epsilon g_u(z)$ for $\epsilon\in\mu_4$.

\begin{lemma}[one place]\label{lem:B3}
In case \textup{(q)} there are a place $u_0\mid3$ of odd residue degree and $\beta\in L_{u_0}$ with
$\Tr_{L_{u_0}/\Q_9}(\beta z)\in\cA$ for every $z\in T(F_{u_0})$.
\end{lemma}

\begin{proof}
Only the places of odd residue degree contribute. For such a place $u$ and $z,z'\in T(F_u)$, varying $y_u$ alone
gives $g_u(z)-g_u(z')\in\cA-\cA\subseteq\frac13O_3$. With $z'=1$ and $z=i$ this says
$(i-1)g_u(1)\in\frac13O_3$, and $i-1$ is a unit at $3$; so $g_u(T(F_u))\subseteq\frac13O_3$. Let $Z_u\subseteq\F_9$
be the set of residues of $3g_u(z)$, $z\in T(F_u)$; it is stable under $\mu_4$, and (q) says that $\sum_uz_u\in A'$
for every choice of $z_u\in Z_u$. Suppose two places $u\ne u'$ had nonzero elements $w\in Z_u$ and $w'\in Z_{u'}$.
Fix $z_{u''}\in Z_{u''}$ at the other places and put $c=\epsilon'w'+\sum z_{u''}$ with $\epsilon'\in\mu_4$. The four
distinct elements $c+\epsilon w$, $\epsilon\in\mu_4$, lie in $A'$, which has four elements; so they are all of $A'$,
and summing them gives $4c=c=\sum A'=0$. As this holds for every $\epsilon'\in\mu_4$, $w'=iw'$, so $w'=0$: a
contradiction. Since $0\notin A'$, exactly one place $u_0$ has $Z_{u_0}\ne\{0\}$; the other $g_u$ take values in
$O_3$, and $\beta=\beta_{u_0}$ works, because $\cA+O_3=\cA$.
\end{proof}

\begin{lemma}[level one]\label{lem:B4}
Let $u\mid3$ have odd residue degree $f$, and let $\beta\in L_u$ satisfy $\Tr_{L_u/\Q_9}(\beta z)\in\cA$ for every
$z\in T(F_u)$. Then $\lambda(x)=3\Tr_{L_u/\Q_9}(\beta x)\bmod3$ is defined on $O_{L_u}$ and vanishes on its maximal
ideal $\mathfrak m$, and the induced $\F_9$-linear map $\lambda_1:\F_{q^2}\to\F_9$ satisfies
$\lambda_1(\mu_{q+1})\subseteq A'$.
\end{lemma}

\begin{proof}
The closed additive span $R$ of $T(F_u)$ is a closed subring of $O_{L_u}$. It contains $\Z_3[\mu_{q+1}]$, which is
the ring of integers $W$ of the unramified extension of $\Q_3$ of degree $2f$ (the order of $3$ modulo $q+1$ is
$2f$). Let $\pi$ be a uniformiser of $F_u$. Then $z=(\pi+i)/(\pi-i)$ and $\sigma(z)=z^{-1}$ lie in $T(F_u)$, and
$2-z-\sigma(z)=4/(\pi^2+1)$ is a unit of $O_{L_u}$ in $R$; its inverse is a limit of its powers, so
$(\pi^2+1)/4\in R$, and $\pi=(z-\sigma(z))(\pi^2+1)/(4i)\in R$. Hence $R\supseteq W[\pi]=O_{L_u}$, since $\pi$ is
also a uniformiser of $L_u$. By continuity $\Tr(\beta O_{L_u})\subseteq\frac13O_3$, so $\lambda:O_{L_u}\to\F_9$ is
defined, $\Z_9$-linear, and zero on $3O_{L_u}=\mathfrak m^e$.

Suppose $\lambda(\mathfrak m)\ne0$, and let $j\ge1$ be the largest integer with $\lambda(\mathfrak m^j)\ne0$.
Identify $\mathfrak m^j/\mathfrak m^{j+1}$ with $\F_{q^2}$ through $x\mapsto(x/\pi^j)\bmod\mathfrak m$; then
$\lambda$ induces a nonzero additive map $\lambda_j:\F_{q^2}\to\F_9$. For $\eta\in O_{F_u}$ the element
$t=(1+i\pi^j\eta)/(1-i\pi^j\eta)$ lies in $T(F_u)$ (as $\sigma(t)=t^{-1}$), and $t\equiv1+2i\pi^j\eta\pmod{\mathfrak
m^{j+1}}$. For $z\in T(F_u)$ with residue $\hat z\in\mu_{q+1}$ this gives $\lambda(zt)=\lambda(z)+\lambda_j(2i\hat z\hat\eta)$, where
$\hat\eta$ is the residue of $\eta$. As $\hat\eta$ runs over $\F_q$, the second term runs over an additive subgroup
$\Sigma_z$ of $\F_9$, and $\lambda(z)+\Sigma_z\subseteq A'$. But $A'=\{a+bi:a,b\ne0\}$ contains no coset of a
nonzero subgroup of $\F_9$, so $\Sigma_z=0$: $\lambda_j$ vanishes on $i\hat z\F_q$ for every $\hat z\in\mu_{q+1}$.
These sets span $\F_{q^2}$ additively, because the $\F_q$-span of $\mu_{q+1}$ is a subfield not contained in
$\F_q$; so $\lambda_j=0$, a contradiction. Hence $\lambda(\mathfrak m)=0$, the induced map $\lambda_1$ is
$\F_9$-linear, and $\lambda_1(\mu_{q+1})\subseteq A'$ because $T(F_u)$ reduces onto $\mu_{q+1}$.
\end{proof}

\begin{lemma}[residue fields]\label{lem:B5}
Let $f$ be odd and $q=3^f$. There is an $\F_9$-linear map $\lambda_1:\F_{q^2}\to\F_9$ with
$\lambda_1(\mu_{q+1})\subseteq A'$ if and only if $f=1$.
\end{lemma}

\begin{proof}
For $f=1$, $\mu_4\subset\F_9$ and $\lambda_1(x)=(1+i)x$ works. Let $f\ge3$. Every $\F_9$-linear map is
$\lambda_1(x)=\Tr_{\F_{q^2}/\F_9}(bx)=\sum_{j<f}(bx)^{9^j}$ for some $b\in\F_{q^2}$. In characteristic $3$,
\[
\lambda_1(z)^4=\lambda_1(z)^3\lambda_1(z)=\sum_{j,k<f}b^{e_{jk}}z^{e_{jk}},\qquad e_{jk}=3\cdot9^j+9^k.
\]
Since $A'=\{y:y^4=-1\}$, if $\lambda_1(\mu_{q+1})\subseteq A'$ then $\Phi(z)=1+\sum_{j,k}b^{e_{jk}}z^{e_{jk}}$
vanishes on $\mu_{q+1}$. On this cyclic group, of order prime to $3$, the characters $z\mapsto z^e$
($e\bmod q+1$) are linearly independent; so for every class $e\bmod q+1$ the sum of the coefficients of the
$z^{e_{jk}}$ with $e_{jk}\equiv e$ vanishes (or equals $-1$, for $e\equiv0$).

Modulo $q+1$, $g=-3$ satisfies $g^f=-3^f\equiv1$ (as $f$ is odd), $9^j\equiv g^{2j}$ and $3\cdot9^j\equiv-g^{2j+1}$.
Let $e_{j'k'}\equiv e_{jj}=4\cdot9^j$. Multiplying by $g^{-2j}$ gives $4\equiv-g^\alpha+g^\gamma$ with
$\alpha\equiv2j'+1-2j$ and $\gamma\equiv2k'-2j\pmod f$, $0\le\alpha,\gamma<f$. As integers, $(-3)^\alpha$ and
$(-3)^\gamma$ have absolute value at most $3^{f-1}$, and $4+2\cdot3^{f-1}<3^f+1$ for $f\ge3$; so the congruence is
the equality $(-3)^\gamma-(-3)^\alpha=4$. It forces $\gamma=0$ (otherwise $3$ divides the left side, or
$(-3)^\gamma=5$), then $(-3)^\alpha=-3$ and $\alpha=1$; so $j'=k'=j$, as $2$ is invertible modulo $f$. Hence the
class of $4\cdot9^j$ contains only $e_{jj}$, it is not $0$ (because $q+1>4$ and $3\nmid q+1$), and its coefficient
$b^{4\cdot9^j}$ must vanish. So $b=0$, and then $\Phi=1\ne0$.
\end{proof}

\begin{proof}[Proof of Theorem~\ref{thm:B}]
If $i\in F$, then (a) and (b) fail ($F(i)=F$, and every residue field above $3$ contains $\F_9$), and
$\chi(F^2)=\infty$ by a theorem of Davies \cite{Davies} (see \cite[Proposition~3]{PP}); so
the theorem holds. Let $i\notin F$.

\emph{Only if.} Let $F^2$ be $3$-colourable. By Lemma~\ref{lem:B2}, (c) or (q) holds. In case (c), Section~\ref{sec:local}
gives (a). In case (q), Lemmas~\ref{lem:B3} and~\ref{lem:B4} give a place $u_0\mid3$ of odd residue degree $f$ and
an $\F_9$-linear $\lambda_1:\F_{q^2}\to\F_9$ with $\lambda_1(\mu_{q+1})\subseteq A'$, and Lemma~\ref{lem:B5} gives
$f=1$: this is (b).

\emph{If.} Under (a) the plane is $2$-colourable \cite[Theorem~A$'$]{MM}, \cite[Theorem~1]{TC}. Under (b), let
$u\mid3$ have residue degree $1$ and let $w$ be the place of $L$ above it. Then $i\notin F_u$, $L_w=F_u(i)$ has
residue field $\F_9$, and the unit vectors of $F_u^2$ are units of $O_w$ whose residues lie in $\mu_4$. A unit step
preserves the class of $z$ modulo $O_w$; fix a representative $\rho$ of each class and colour $z$ by
$\ell((z-\rho)\bmod\mathfrak m_w)$, where $\ell(a+bi)=a+b$ on $\F_9=\F_3+\F_3i$. Since $\ell(\pm1)=\pm1$ and
$\ell(\pm i)=\pm1$, every unit step changes the colour. This colours $F_u^2\supseteq F^2$ with three colours.
Finally $\chi(F^2)\le2$ if and only if (a), by \cite[Theorem~1]{TC}.
\end{proof}

The proof uses Theorem~W$^+$ only through Lemma~\ref{lem:B1}, for the finite sets $G_NV$; the rest is
compactness, weak approximation for the norm-one torus, the arithmetic of the places above $2$ and $3$, and, in case
(c), the two-colour theorem of \cite{TC}.

\section{Consequences}\label{sec:cor}

\begin{corollary}[no circular chromatic number between $2$ and $3$]\label{cor:B6}
Let $F$ be a number field. If $F^2$ has a homomorphism to a circular clique $K_{p/q}$ with $p/q<3$, then
$\chi(F^2)\le2$. Hence $\chi_c(F^2)=2$, $\chi_c(F^2)=3$ or $\chi_c(F^2)>3$, according as $\chi(F^2)=2$, $3$ or at
least $4$; and $F^2$ has a homomorphism to a cycle of odd length at least $5$ only if it is bipartite.
\end{corollary}

\begin{proof}
If $i\in F$ there is nothing to prove, so let $i\notin F$. For every finite $U$ as in the proof of
Lemma~\ref{lem:B1}, Theorem~W$^+$ gives a character with $\xi(U)\subseteq[q/p,1-q/p]+\Z$, and
$q/p>1/3$. In that proof $\varphi(v)$ then satisfies $\operatorname{Re}(\overline{\varphi(v)})\in[q/p,1-q/p]+\Z$ (take
$\gamma=1$), while a point $N\varepsilon$ with $\varepsilon\in E_q$ (by Proposition~\ref{prop:shape} these are the
points of type~q) has real part in $\{\frac13,\frac23\}+\Z$, as $3\nmid N$. So every $\varepsilon_v$ has type~c, the conditions $P_v$ in the proof of Lemma~\ref{lem:B2} can all
be taken with $\tau(v)=\mathrm c$, and (c) holds; so $\chi(F^2)\le2$. For the second statement: if
$\chi_c(F^2)<3$, then by the definition of the infimum $F^2$ maps to some $K_{p/q}$ with $p/q<3$, so $\chi(F^2)\le2$;
hence $\chi_c(F^2)=3$ when $\chi(F^2)=3$, as $\chi_c\le\chi$. If $\chi(F^2)\ge4$, then by the theorem of de Bruijn and
Erd\H{o}s \cite{dBE} some finite subgraph $H$ has $\chi(H)\ge4$, so $\chi_c(F^2)\ge\chi_c(H)>\chi(H)-1\ge3$.
\end{proof}

Call a number field $F$ \emph{split above $2$ and $3$} if $i\in F_v$ for every place $v$ of $F$ above $2$ and above
$3$, that is, if every such place splits in $F(i)$; then (a) and (b) fail. The admissible fields of
\cite[Section~4]{TC}, which must moreover satisfy $\sqrt3\notin F$ and $i\in F_v$ at the places above $7$, $11$ and
$19$, are split above $2$ and $3$.

\begin{corollary}\label{cor:B7}
Every number field $F$ that is split above $2$ and $3$ has $\chi(F^2)\ge4$.
\end{corollary}

A real quadratic field $\Q(\sqrt d)$ is split above $2$ and $3$ exactly when $d\equiv23\pmod{24}$, so for quadratic
fields Corollary~\ref{cor:B7} is contained in \cite{FC}. Of the biquadratic fields listed as admissible in the account
of \cite[\nolinkurl{notes/local_global.md}, Section~4]{Repo}, $\Q(\sqrt2,\sqrt{47})$, $\Q(\sqrt{11},\sqrt{13})$ and
$\Q(\sqrt{10},\sqrt{38})$ (which contains $\sqrt{95}$) contain some $\sqrt d$ with $d\equiv11\pmod{12}$; for
$\Q(\sqrt2,\sqrt{31})$ and $\Q(\sqrt2,\sqrt{55})$ the bound is new.

\begin{corollary}[local fields]\label{cor:local}
Let $K$ be a finite extension of $\Q_p$. Then $\chi(K^2)=2$ if $p=2$ and $K(i)/K$ is ramified, $\chi(K^2)=3$ if
$p=3$ and the residue degree of $K$ is $1$, and $\chi(K^2)\ge4$ in every other case.
\end{corollary}

\begin{proof}
In the first case the $2$-colouring of the proof of \cite[Proposition~3]{TC} colours $K^2$ itself (that proof
colours $F_v^2$, and it applies verbatim to every finite extension $K$ of $\Q_2$ with $K(i)/K$ ramified); in the
second case the $3$-colouring of the proof of Theorem~\ref{thm:B} colours $K^2$ itself; and
$\chi(K^2)\ge3$ when $p=3$, because $\Q(\sqrt7)\subset\Q_3\subseteq K$ and $\chi(\Q(\sqrt7)^2)=3$. If $p\ge5$, then
$\chi(K^2)\ge\chi(\Q_p^2)\ge4$ by \cite[Corollary~3]{FC}. If $i\in K$, then $\chi(K^2)=\infty$ \cite[Proposition~3]{PP}.
There remain the case $p=2$, $i\notin K$, $K(i)/K$ unramified, and the case $p=3$ with odd residue degree at least
$3$. In both we find a number field $F\subset K$ for which (a) and (b) fail; then $\chi(K^2)\ge\chi(F^2)\ge4$ by
Theorem~\ref{thm:B}.

Let $m=[K:\Q_p]$ and let $p'$ be the other element of $\{2,3\}$. If $p=2$, then $m$ is even: if $m$ were odd,
$\Q_2^\times/\Q_2^{\times2}$ would inject into $K^\times/K^{\times2}$, and since $-1$ and $-5$ are not squares in
$\Q_2$, $i\notin K$ and $K(i)\ne K(\sqrt5)$, which is the unramified quadratic extension of $K$; so $K(i)/K$ would be
ramified. Choose monic polynomials over $\Q_p$, $\Q_{p'}$ and $\Q_5$ of a common even degree $n$:
\begin{itemize}
\item over $\Q_p$, $h_Kh'$, where $h_K$ is the minimal polynomial of a generator of $K$ and $h'$ is a product of
distinct irreducible quadratics, coprime to $h_K$, whose roots generate $\Q_2(i)$ (if $p=2$) or $\Q_9$ (if $p=3$);
if $p=3$ and $m$ is odd, $h'$ also contains one cubic factor, coprime to $h_K$, generating the unramified cubic
extension $\Q_{27}$ of $\Q_3$;
\item over $\Q_{p'}$, a product of distinct irreducible quadratics whose roots generate $\Q_2(i)$, resp.\ $\Q_9$;
\item over $\Q_5$, an irreducible polynomial.
\end{itemize}
By weak approximation choose a monic $h\in\Q[x]$ whose coefficients are so close to the three targets that, by
Krasner's lemma, $h$ factors over $\Q_p$, $\Q_{p'}$ and $\Q_5$ with the same degrees and the same fields as the
targets. Then $F=\Q[x]/(h)$ is a field (as $h$ is irreducible over $\Q_5$); it embeds in $K$, through a root of
$h$ close to the chosen generator of $K$; its places above $p$ are the place that gives $K$ and places that contain
$i$ (or, for the cubic factor, have residue degree $3$); and its places above $p'$ all contain $i$. So no place
above $2$ ramifies in $F(i)$ and no place above $3$ has residue degree $1$.
\end{proof}

For example $\chi(\Q_{27}^2)\ge4$, although $i\notin\Q_{27}$ and $x^2+y^2$ is anisotropic over $\Q_{27}$, as over
$\Q_3$; and $\chi(\Q_2(\sqrt3)^2)\ge4$, as $\Q_2(\sqrt3,i)=\Q_2(\sqrt3,\sqrt{-3})$ is unramified over
$\Q_2(\sqrt3)$ (this also follows from $\chi(\Q(\sqrt{11})^2)=4$, since $\Q_2(\sqrt{11})=\Q_2(\sqrt3)$).
Corollary~\ref{cor:local} extends \cite[Corollary~3]{FC} ($p\ge5$) and the values $\chi(\Q_2^2)=2$ and
$\chi(\Q_3^2)=3$, which follow from Madore's results \cite{Madore} (see \cite[Theorem~2]{PP}).

\section{An explicit family: two square roots}\label{sec:two}

For many fields the general argument can be replaced by two explicit relations. For $c\ge1$ put
\[
u_c=\frac{(1+i\sqrt c)^2}{1+c}=\frac{(1-c)+2\sqrt c\,i}{1+c},\qquad\text{so that}\qquad(1+c)(u_c+\bar u_c)=2(1-c).
\]

\begin{proposition}\label{prop:two}
Let $F\subset\R$ contain $\sqrt a$ and $\sqrt b$, where $a\equiv2\pmod3$ and $b\equiv7\pmod8$ are positive integers
\textup{(}$a=b$ is allowed\textup{)}, and let $N=5^k>\frac{4\sqrt{10}}3\max(a,b)$. Then no character of the group
generated by $U=G_N\{1,u_a,\bar u_a,u_b,\bar u_b\}$ maps every element of $U$ into $[1/3,2/3]+\Z$. Hence
$\chi(F^2)\ge4$.
\end{proposition}

\begin{proof}
As in the proof of Lemma~\ref{lem:B1}, such a character gives numbers $w_v=\varphi(v)\in S_N$
($v\in\{1,u_a,\bar u_a,u_b,\bar u_b\}$) that satisfy the two relations exactly. The sums of the absolute values of
their coefficients are $4a$ and $4b$, and $6r\cdot4c=\frac{4\sqrt{10}}3c<N$; so Lemma~\ref{lem:exact} gives, for the
types,
\[
(1+a)(\varepsilon_{u_a}+\varepsilon_{\bar u_a})=2(1-a)\varepsilon_1,\qquad(1+b)(\varepsilon_{u_b}+\varepsilon_{\bar u_b})=2(1-b)\varepsilon_1.
\]
\emph{At $3$.} $3\mid1+a$ and $3\nmid2(1-a)$. By \eqref{eq:types} the left side of the first relation has
$v_3\ge0$, hence $v_3(\varepsilon_1)\ge0$, and $\varepsilon_1$ is not of type~q.
\emph{At $2$.} $v_2(1+b)\ge3$ and $v_2(2(1-b))=2$, as $1-b\equiv2\pmod8$. By \eqref{eq:types} the left side of the
second relation has $v_{1+i}\ge2\cdot3-1=5$, while $v_{1+i}(2(1-b)\varepsilon_1)=4+v_{1+i}(\varepsilon_1)=3$ if
$\varepsilon_1$ has type~c. So $\varepsilon_1$ is not of type~c either, a contradiction.
\end{proof}

For $a=b=d$ this is the case $d\equiv23\pmod{24}$ of \cite[Theorem~1]{FC}, where \cite[Theorem~4(1)]{FC} needs
only $d<21\cdot5^{k-1}$. Proposition~\ref{prop:two} covers, for
example, $\Q(\sqrt2,\sqrt7)$, $\Q(\sqrt2,\sqrt{31})$, $\Q(\sqrt5,\sqrt7)$, $\Q(\sqrt2,\sqrt{55})$ and
$\Q(\sqrt3,\sqrt5)$ (with $b=15$), and every field containing one of them. In terms of Theorem~\ref{thm:B}, $\sqrt a$
makes every place above $3$ have even residue degree, and $\sqrt b$ puts $i$ in every completion above $2$.

\section{Circular colourings}\label{sec:circ}

For the plane over $\R$, DeVos, Ebrahimi, Ghebleh, Goddyn, Mohar and Naserasr proved $\chi_c(\R^2)\ge4$
\cite{DeVos}; since de Grey's theorem gives a finite subgraph with $\chi=5$, in fact $\chi_c(\R^2)>4$. For planes over
number fields the values $2$ and $3$ are immediate (a bipartite plane; a plane with $\chi=3$ containing a unit
triangle), and by Corollary~\ref{cor:B6} the planes with $\chi_c\le3$ are exactly those with $\chi\le3$. We know of no
earlier non-integral value. Theorem~\ref{thm:C} shows that $7/2$ occurs; its upper bound comes from a residue field.

\begin{proposition}\label{prop:C1}
Let $v$ be a place of $F$ whose residue field is $\F_p$, with $p\equiv3\pmod4$, and let $\lambda:\F_{p^2}\to\F_p$ be
$\F_p$-linear with $\lambda(\mu_{p+1})\subseteq\{m,m+1,\dots,p-m\}$. Then $F_v^2$, hence $F^2$, has a homomorphism
to $K_{p/m}$. For $p=7$, $\lambda(a+bi)=2a+3b$ takes the values $2,3,4,5$ on $\mu_8\subset\F_{49}$, so $F^2$ maps to
$K_{7/2}$, and $\chi_c(F^2)\le7/2$.
\end{proposition}

\begin{proof}
As $-1$ is not a square in $\F_p$, $i\notin F_v$, and $L_w=F_v(i)$ is unramified over $F_v$ with residue field
$\F_{p^2}$. A unit vector $z$ has $z\bar z=1$, so it is a unit of $O_w$, and since the conjugation of $L_w/F_v$ reduces
to the Frobenius $x\mapsto x^p$, its residue satisfies $\hat z^{p+1}=1$. A unit step preserves the class modulo
$O_w$; fix a representative $\rho$ of each class and colour $z$ by $\lambda((z-\rho)\bmod\mathfrak m_w)\in\Z/p$. A unit
step changes the colour by an element of $\{m,\dots,p-m\}$. For $p=7$, $\mu_8=\{\pm1,\pm i,\pm2\pm2i\}$, and $2a+3b$
takes the values $2,5,3,4,3,5,2,4$ at $1,-1,i,-i,2+2i,2-2i,-2+2i,-2-2i$.
\end{proof}

\begin{proof}[Proof of Theorem~\ref{thm:C}]
\emph{Upper bounds.} $7$ splits in $\Q(\sqrt{11})$ ($11\equiv2^2\pmod7$) and ramifies in $\Q(\sqrt{35})$, so both
fields, and $\Q_7$, have a place with residue field $\F_7$; Proposition~\ref{prop:C1} applies.

\emph{Lower bounds.} For $d\in\{11,35\}$ let $u_n=(n+i\sqrt d)^2/(n^2+d)$, $V=\{1,u_1,\bar u_1,u_7,\bar u_7,u_{19},\bar
u_{19}\}$ and $U=G_{25}V$, a set of $140$ unit vectors ($70$ up to sign). We checked by computer that no character of
$\Z U$ maps every element of $U$ into the open interval $(2/7,5/7)+\Z$. The certificate branches on the integer
values of the relations among the vectors ($66$ relations; $17\,744$ nodes for $d=11$ and $14\,199$ for $d=35$) and
closes every leaf with an exact Farkas vector; two exact checkers that share no code accept it
(Section~\ref{sec:checks}). If $F^2$, $F=\Q(\sqrt d)$, mapped to $K_{p/q}$ with $p/q<7/2$, Theorem~W$^+$ would
give a character with values in $[q/p,1-q/p]\subset(2/7,5/7)$ on $U$. So $\chi_c(F^2)\ge7/2$. Since
$\Q(\sqrt{11})\subset\Q_7$, also $\chi_c(\Q_7^2)\ge7/2$.

\emph{Chromatic numbers.} $K_{7/2}\to K_4$ gives $\chi\le4$, and $\chi\ge\chi_c=7/2>3$ gives $\chi\ge4$ (for the two
quadratic fields this is also \cite[Theorem~1]{FC}, for $\Q_7$ \cite[Theorem~2]{PP}). The last sentence of the
theorem is Corollary~\ref{cor:B6}.
\end{proof}

The lower bound is about the infinite graph $\Cay(\Z U,U)$, to which Theorem~W$^+$ is applied directly; we do
not know a finite unit-distance graph with circular chromatic number exactly $7/2$. The open interval is essential:
the characters that come from the $7$-adic colourings take values in the closed interval $[2/7,5/7]$ on $U$. The
certificates were found with a floating-point mixed-integer solver, which reports
$\max_\xi\min_{u\in U}\|\xi(u)\|=2/7$ for both sets; nothing in the proof depends on it.

The circular chromatic number separates planes with the same chromatic number:

\begin{proposition}\label{prop:fiftynine}
$\chi_c(\Q(\sqrt{59})^2)\ge53/15>7/2$, while $\chi(\Q(\sqrt{59})^2)=4$.
\end{proposition}

\begin{proof}
For $d=59$ the same set $U=G_{25}V$ has no character into the open interval $(15/53,38/53)$: the certificate has
$127\,181$ nodes and $105\,761$ leaves, and both checkers accept it (Section~\ref{sec:checks}). As in the proof of
Theorem~\ref{thm:C}, Theorem~W$^+$ gives $\chi_c\ge53/15$. Here $7$ is inert, so Proposition~\ref{prop:C1} does not
apply. We have $\chi\ge4$ by \cite[Theorem~1]{FC}, as $59\equiv11\pmod{12}$. And $\chi\le4$: the place of
$F=\Q(\sqrt{59})$ above $2$ has residue field $\F_2$ and is inert in $F(i)$ (as $59\equiv3\pmod8$), so the unit
vectors reduce into the group $\mu_3$ of elements of norm one in $\F_4$, and, as in the proof of
Proposition~\ref{prop:C1}, the residue modulo the maximal ideal is a proper $4$-colouring (a bound of Fischer
\cite{Fischer1990}).
\end{proof}

\begin{corollary}\label{cor:sevenhalves}
$\chi_c(E^2)=7/2$ for every finite extension $E$ of $\Q_7$ with residue degree $1$, and $\chi_c(F^2)=7/2$ for every
number field $F$ that contains $\sqrt{11}$ or $\sqrt{35}$ and has a place above $7$ of residue degree $1$, for example
$\Q(\sqrt7,\sqrt{11})$ and $\Q(\sqrt2,\sqrt{35})$.
\end{corollary}

\begin{proof}
$E\supseteq\Q_7\supset\Q(\sqrt{11})$ and $F\supseteq\Q(\sqrt{11})$ or $\Q(\sqrt{35})$ give $\chi_c\ge7/2$ by
Theorem~\ref{thm:C}; the residue field $\F_7$ gives $\chi_c\le7/2$ by Proposition~\ref{prop:C1} (whose proof
applies to $E$ itself).
\end{proof}

\subsection*{Locally constant colourings}
The values $2$, $3$ and $7/2$ are the only ones below $4$ that colourings through a single completion can give
(Corollary~\ref{cor:localvalues}). Let
$E$ be a finite extension of $\Q_p$ with $i\notin E$ and $E(i)/E$ unramified (for $p$ odd this always holds; for $p=2$
it is the case where (a) fails), let $\F_q$ be the residue field of $E$ ($q=p^f$), $w$ the place of $L=E(i)$, and
\[
\kappa_1(E)=\max_\lambda\min_{z\in\mu_{q+1}}\|\lambda(z)\|,
\]
the maximum over the additive maps $\lambda:\F_{q^2}\to\frac1p\Z/\Z$.

\begin{proposition}\label{prop:levels}
Suppose $E^2$ has a locally constant homomorphism to $K_{P/Q}$, $P/Q<4$. Then $Q/P\le\kappa_1(E)$. Conversely $E^2$ maps to $K_{p/m}$ when $m/p=\kappa_1(E)>0$.
\end{proposition}

\begin{proof}
The converse is the colouring of Proposition~\ref{prop:C1}, with $\F_q$ in place of $\F_p$. Let $\pi$ be a
uniformiser of $E$ (also one of $L$) and $\sigma$ the non-trivial automorphism of $L/E$. On the compact set
$O_w$ the homomorphism is constant on the cosets of some power $\mathfrak m^k$ of the maximal ideal, and only its
restriction to $O_w$ is used. It induces a homomorphism from $\Cay(O_w/\mathfrak m^k,\bar T)$, where
$\bar T$ is the image of $T(E)$, to $K_{P/Q}$; by Theorem~W$^+$ there is a character $\xi$ of
$O_w/\mathfrak m^k$ with $\|\xi(z)\|\ge Q/P>1/4\ge1/(2p)$ for every $z\in T(E)$. We show that $\xi$ vanishes on
$\mathfrak m$; then $\xi$ is a map $\lambda$ as in the definition of $\kappa_1$, and $Q/P\le\kappa_1(E)$, since $T(E)$
reduces onto $\mu_{q+1}$ (it contains the Teichm\"uller lifts of $\mu_{q+1}$, on which $\sigma$ acts as
$x\mapsto x^q$). Suppose not, and let $J\ge1$ be the largest integer with $\xi(\mathfrak m^J)\ne0$. Through
$x\mapsto(x/\pi^J)\bmod\mathfrak m$, $\xi$ induces a nonzero additive map $\lambda_J:\F_{q^2}\to\frac1p\Z/\Z$. For
$x\in O_w$, $t=(1+\pi^Jx)/\sigma(1+\pi^Jx)$ lies in $T(E)$ and $t\equiv1+\pi^J(x-\sigma(x))\pmod{\mathfrak m^{J+1}}$,
and the residues of $x-\sigma(x)$ fill the $\F_q$-line $D=\{y-y^q:y\in\F_{q^2}\}$. For $z\in T(E)$ with residue
$\hat z$, $\xi(zt)=\xi(z)+\lambda_J(\hat zd)$, $d\in D$; the values $\lambda_J(\hat zD)$ form a subgroup of
$\frac1p\Z/\Z$, and if it were the whole group, $\xi(T(E))$ would contain a coset of $\frac1p\Z/\Z$ and a point
with $\|\cdot\|\le\frac1{2p}$. So $\lambda_J$ vanishes on $\hat zD$ for every $\hat z\in\mu_{q+1}$, hence on their
additive span, which is $D$ times the $\F_q$-span of $\mu_{q+1}$, that is $\F_{q^2}$: a contradiction.
\end{proof}

\begin{proposition}\label{prop:residue}
$\kappa_1(E)>1/4$ only when $p=3$ and $f=1$, where $\kappa_1=1/3$, or $p=7$ and $f=1$, where $\kappa_1=2/7$.
In every other case with $i\notin E$ and $E(i)/E$ unramified, $\kappa_1(E)<1/4$; it is $0$ when $p=2$ or $f\ge3$.
\end{proposition}

\begin{proof}
For $p=2$: if $\lambda(\mu_{q+1})=\{\frac12\}$, summing over $\mu_{q+1}$, whose elements add up to $0$, gives
$\lambda(0)=\frac{q+1}2\equiv\frac12$, as $q+1$ is odd; so some $z$ has $\lambda(z)=0$.

Let $p$ be odd; then $p\equiv3\pmod4$ and $f$ is odd, as $i\notin E$. Let $\psi(x)=e^{2\pi ix/p}$ on $\F_p$, and
$S(b)=\sum_{z\in\mu_{q+1}}\psi(\Tr_{\F_{q^2}/\F_p}(bz))$ for $b\in\F_{q^2}^\times$. Every additive
$\lambda:\F_{q^2}\to\frac1p\Z/\Z$ is $z\mapsto\frac1p\Tr(bz)$ for some $b$. Writing the indicator function of
$\mu_{q+1}=\ker N_{\F_{q^2}/\F_q}$ with the multiplicative characters $\chi$ of $\F_q^\times$, and using the
Davenport--Hasse relation $G(\chi\circ N,\psi\circ\Tr)=-G(\chi,\psi_q)^2$ for Gauss sums \cite{DH}
($\psi_q=\psi\circ\Tr_{\F_q/\F_p}$; see \cite[Chapter~5]{LN}), one finds
\[
S(b)=\frac1{q-1}\sum_\chi\bar\chi(N(b))\,G(\chi\circ N,\psi\circ\Tr)=-\sum_{y\in\F_q^\times}\psi_q\Bigl(y+\frac{N(b)}y\Bigr),
\]
minus a Kloosterman sum; so $|S(b)|\le2\sqrt q$ by Weil's bound \cite{Weil}. Hence, for $b\ne0$, the number of $z\in\mu_{q+1}$
with $\Tr(bz)=0$ is $\frac1p\sum_{a\in\F_p}\sum_z\psi(a\Tr(bz))=\frac{q+1}p+\frac1p\sum_{a\ne0}S(ab)\ge\frac{q+1}p-
\frac{p-1}p\,2\sqrt q$, which is positive when $f\ge3$ (then $\sqrt q\ge p^{3/2}>2(p-1)$); so $\kappa_1=0$ for $f\ge3$.

Let $f=1$, $p=4k+3$, and $I=\{x:|x|\le k\}\subset\F_p$, the residues with $\|x/p\|<\frac14$. The number of
$z\in\mu_{p+1}$ with $\Tr(bz)\in I$ is
\[
\frac{(p+1)|I|}p+\frac1p\sum_{a\ne0}\Bigl(\sum_{x\in I}\psi(-ax)\Bigr)S(ab)
\ \ge\ \frac{p^2-1}{2p}-\frac{2\sqrt p}p\sum_{a=1}^{p-1}\frac1{|\sin(\pi a/p)|}
\ \ge\ \frac{p^2-1}{2p}-2\sqrt p\,(\ln p+1),
\]
because $|\sum_{x\in I}\psi(-ax)|=|\sin(\pi a(2k+1)/p)/\sin(\pi a/p)|\le1/|\sin(\pi a/p)|$, and, by the symmetry
$a\leftrightarrow p-a$ and $\sin x\ge2x/\pi$ on $[0,\pi/2]$, $\sum_{a=1}^{p-1}1/|\sin(\pi a/p)|\le
p\sum_{a\le(p-1)/2}1/a\le p(\ln p+1)$. The last expression is positive for $p=1001$ and increasing for $p\ge1001$
(its derivative is $\frac12+\frac1{2p^2}-\frac{\ln p+3}{\sqrt p}>0$ there). For the primes $p\equiv3\pmod4$ with
$11\le p<3000$ a direct computation
(in $\F_{p^2}=\F_p(i)$, where $bz$ runs over the circle $s^2+t^2=N(b)$ and $\Tr(bz)=2s$) gives $\kappa_1<\frac14$;
the largest value is $\frac4{19}$. For $p=3$ and $p=7$ a direct computation gives $\kappa_1=\frac13$ and
$\kappa_1=\frac27$, attained by the maps of the proof of Theorem~\ref{thm:B} and of Proposition~\ref{prop:C1}.
\end{proof}

\begin{corollary}[local circular values]\label{cor:localvalues}
Let $E$ be a finite extension of $\Q_p$. If $E^2$ has a locally constant homomorphism to $K_{P/Q}$ with $P/Q<4$, then
either $p=2$ and $E(i)/E$ is ramified, or $(p,f)=(3,1)$ and $P/Q\ge3$, or $(p,f)=(7,1)$ and $P/Q\ge7/2$.
\end{corollary}

\begin{proof}
If $i\in E$, then $\chi(E^2)=\infty$ \cite[Proposition~3]{PP}. Otherwise combine Propositions~\ref{prop:levels}
and~\ref{prop:residue}.
\end{proof}

So the restriction to $F^2$ of a locally constant colouring of a completion $F_v^2$ of a number field gives, below
$4$, only the bounds $2$ (condition (a)), $3$ (condition (b)) and $7/2$ (a place above $7$ of residue degree $1$).

\section{A gap above \texorpdfstring{$3$}{3}}\label{sec:gap}

By Corollary~\ref{cor:B6}, $\chi_c(F^2)$ is $2$, $3$ or larger than $3$. Theorem~\ref{thm:D} says that in the last
case it stays away from $3$. Its proof is that of Theorem~\ref{thm:B} with $I=[1/3,2/3]$ replaced by a slightly
wider interval; the new ingredient is the shape of the characters along the probe for such intervals.

\emph{Notation.} In this section $N=5^k$ with $k\ge1$, $\rho=(3+4i)/5$, $h=(1+i)/2$ and
$\Lambda_\pm=\frac{2\pm i}5\Z[i]$. For $\frac14<\theta\le\frac13$ let $S_N^\theta$ be defined as $S_N$ with
$[\theta,1-\theta]$ in place of $I$, and put
\[
s=\tfrac12-\theta,\qquad\delta=\tfrac13-\theta,\qquad B_s=\{a+bi:\ a,b\in\R,\ |a|,|b|\le s\},
\]
so that $s=\frac16+\delta$, and $P_k=\{x\in\C:\ \bar x\rho^j\in B_s\text{ for }|j|\le k\}$. For $\varepsilon\in E_q$
and $|j|\le k$, Lemma~\ref{lem:values} below gives $\eta^N_j(\varepsilon)\in\{\pm1\pm i\}$ with
$\overline{N\varepsilon}\rho^j\in h+\eta^N_j(\varepsilon)/6+\Z[i]$; for $1\le m\le k$ let
\[
Y^N_m(\varepsilon)=\{y\in\C:\ \eta^N_j(\varepsilon)/6+\bar y\rho^j\in B_s\text{ for }|j|\le m\},
\]
and $Y_k(\varepsilon)=Y^N_k(\varepsilon)$. These depend on $\theta$, on $N$, and on $\varepsilon$ only modulo
$\Z[i]$; when $N$ is fixed we write $\eta_j(\varepsilon)$ for $\eta^N_j(\varepsilon)$. We write $w_1=\operatorname{Re}w$
and $w_2=\operatorname{Im}w$.

\begin{lemma}\label{lem:values}
\textup{(i)} $G_N=\{i^a\rho^j:a\in\Z,\ |j|\le k\}$, and $c\in S_N^\theta$ if and only if $\bar c\rho^j\in
h+B_s+\Z[i]$ for every $|j|\le k$.

\textup{(ii)} If $\varepsilon\in E_c$ and $|j|\le k$, then $\overline{N\varepsilon}\rho^j\in h+\Z[i]$.

\textup{(iii)} If $\varepsilon=\frac{\alpha+\beta i}3+t\in E_q$ and $|j|\le k$, then
$\overline{N\varepsilon}\rho^j\equiv(-1)^k(-i)^j\frac{\alpha-\beta i}3\pmod{\Z[i]}$. This lies in
$h+\eta/6+\Z[i]$ for a unique $\eta\in\{\pm1\pm i\}$; moreover $\eta_{j-1}(\varepsilon)=i\,\eta_j(\varepsilon)$ for
$1-k\le j\le k$ (so $\eta_{j-2}(\varepsilon)=-\eta_j(\varepsilon)$), and $\eta_j(\bar\varepsilon)=
\overline{\eta_{-j}(\varepsilon)}$.

\textup{(iv)} $N\varepsilon+P_k\subseteq S_N^\theta$ for $\varepsilon\in E_c$, and $N\varepsilon+Y_k(\varepsilon)\subseteq
S_N^\theta$ for $\varepsilon\in E_q$.
\end{lemma}

\begin{proof}
(i) If $\gamma\in G_N$, then $N\gamma\in\Z[i]$ has norm $5^{2k}$, so $N\gamma=u(2+i)^a(2-i)^b$ with a unit $u$ and
$a+b=2k$, and $\gamma=u\big(\frac{2+i}{2-i}\big)^{a-k}=u\rho^{a-k}$; conversely
$N\rho^{\pm j}=5^{k-j}(3\pm4i)^j\in\Z[i]$ for $0\le j\le k$. Since $\operatorname{Re}(\pm iw)=\mp\operatorname{Im}w$
and $[\theta,1-\theta]+\Z=\frac12+[-s,s]+\Z$ is symmetric, $\operatorname{Re}(\bar c\gamma)\in[\theta,1-\theta]+\Z$
for all $\gamma\in G_N$ if and only if both coordinates of every $\bar c\rho^j$ lie in $\frac12+[-s,s]+\Z$.

(ii) Write $\varepsilon=h+t$. The Gaussian integer $N\rho^j$ has odd norm, so $N\rho^j=1+(1+i)m$ with
$m\in\Z[i]$, and $\bar hN\rho^j=\bar h+m\equiv h$, as $\bar h=h-i$; and $\bar tN\rho^j\in\Z[i]$.

(iii) $\overline{N\varepsilon}\rho^j\equiv\frac{\alpha-\beta i}3N\rho^j$, and $N\rho^j\equiv(-1)^k(-i)^j\pmod3$,
because $5\equiv-1$ and $3\pm4i\equiv\pm i\pmod3$. Modulo $\Z[i]$, $\frac{\alpha-\beta i}3$ times a unit has both
coordinates in $\{\frac13,\frac23\}$, so it lies in $h+\frac{\pm1\pm i}6+\Z[i]$. Since $(-i)^{j-1}=i(-i)^j$ and
$ih\equiv h$, $\eta_{j-1}=i\eta_j$. Finally $\overline{N\bar\varepsilon}\rho^j$ is the conjugate of
$\overline{N\varepsilon}\rho^{-j}$, and $\bar h\equiv h$.

(iv) follows from (i)--(iii).
\end{proof}

\begin{lemma}\label{lem:radii}
\textup{(i)} $P_1=sP$, where $P$ is the $12$-gon with vertices $\pm1\pm\frac i3$, $\pm\frac13\pm i$ and
$\frac57(\pm1\pm i)$; so $|x|\le s\sqrt{10}/3$ for $x\in P_k$.

\textup{(ii)} $|y|\le\sqrt2\,\delta$ for $y\in Y^N_m(\varepsilon)$, $1\le m\le k$.
\end{lemma}

\begin{proof}
(i) $P_1$ is defined by the twelve inequalities $\pm(\bar x\rho^j)_e\le s$, $j\in\{0,\pm1\}$, $e\in\{1,2\}$, which
are homogeneous in $(x,s)$; intersecting their boundary lines in pairs gives the vertices. For instance at
$x=1+\frac i3$ we have $\bar x=1-\frac i3$, $\bar x\rho=\frac{13}{15}+\frac35i$ and $\bar x\rho^{-1}=\frac13-i$.
Every vertex has modulus at most $\sqrt{10}/3$, and $P_k\subseteq P_1$.

(ii) As $s=\frac16+\delta$, the condition $\eta_j/6+\bar y\rho^j\in B_s$ implies
$\eta_{j,e}(\bar y\rho^j)_e\le\delta$ for $e=1,2$, where $\eta_j=\eta_{j,1}+\eta_{j,2}i$. So
$Y^N_m(\varepsilon)\subseteq Y^N_1(\varepsilon)\subseteq\delta X(\varepsilon)$, where $X(\varepsilon)$ is defined by
the six inequalities $\eta^N_{j,e}(\bar y\rho^j)_e\le1$, $|j|\le1$ (for $\delta=0$ read $\delta X(\varepsilon)$ as
$\{0\}$: as $X(\varepsilon)$ is bounded, the same inequalities with right-hand side $0$ have only the solution $0$).
By Lemma~\ref{lem:values}(iii), $\eta^N_j(\varepsilon)=\eta^5_j(\pm\varepsilon)$ for $|j|\le1$; replacing
$\varepsilon$ by $i\varepsilon$ replaces $\eta_j$ by $-i\eta_j$ and $X$ by $iX$, and multiplication by $i$ permutes
the four classes of $E_q$ modulo $\Z[i]$ cyclically. So every $X(\varepsilon)$ is a rotation of $X(\frac{1+i}3)$ at
level $5$, where $\eta_{-1}=1+i$, $\eta_0=1-i$, $\eta_1=-1-i$. This is the hexagon with vertices $-\frac54$, $-\frac54i$,
$-\frac57(1+i)$, $1-\frac i2$, $1+i$ and $-\frac12+i$, of moduli at most $\sqrt2$.
\end{proof}

\begin{lemma}\label{lem:gapC}
Let $0<s<11/56$ and $k\ge1$, and let $x\in\C$ with $|x|\le s\sqrt2$ and $\bar x\rho^{k-2}\in B_s$. Then
$\bar x\rho^k+\nu\notin B_s+\Z[i]$ for every $\nu\in\Lambda_+\setminus\Z[i]$. The same holds with $\rho^{-k}$,
$\rho^{2-k}$ and $\Lambda_-$ in place of $\rho^k$, $\rho^{k-2}$ and $\Lambda_+$.
\end{lemma}

\begin{proof}
Let $u=\bar x\rho^k$ and $w=\bar x\rho^{k-2}=u\rho^{-2}=u(-7-24i)/25$, so that $w_1=(-7u_1+24u_2)/25$ and
$w_2=-(24u_1+7u_2)/25$. Suppose $u+\nu\in B_s+\Z[i]$; changing $\nu$ by a Gaussian integer we may assume
$u+\nu\in B_s$. The non-zero classes of $\Lambda_+/\Z[i]$ are represented by $\pm\frac{2+i}5$ and
$\pm\frac{-1+2i}5$, whose coordinates have absolute value at most $\frac25$; every other element of the class has a
coordinate of absolute value at least $\frac35$, which $u_e+\nu_e\in[-s,s]$ excludes, as $|u_e|\le|x|\le s\sqrt2$ and
$s(1+\sqrt2)<\frac35$. If $\nu=\pm\frac{2+i}5$, then $\mp u_1\ge\frac25-s$ and $\mp u_2\ge\frac15-s$, so
$\pm w_2=\mp(24u_1+7u_2)/25\ge(11-31s)/25$. If $\nu=\pm\frac{-1+2i}5$, then $\pm u_1\ge\frac15-s$ and $\mp
u_2\ge\frac25-s$, so $\mp w_1=\pm(7u_1-24u_2)/25\ge(11-31s)/25$. In both cases
$|w_e|\ge(11-31s)/25>s$, as $56s<11$, contradicting $w\in B_s$. The second statement follows by complex conjugation,
which exchanges $\rho$ and $\rho^{-1}$, and $\Lambda_+$ and $\Lambda_-$, and preserves $B_s$ and $\Z[i]$ (apply the
first to $\bar x$ and $\bar\nu$).
\end{proof}

\begin{lemma}\label{lem:gapQ}
Let $0\le\delta<5/168$, $k\ge2$ and $\varepsilon\in E_q$, and let $y\in\C$ with $|y|<\frac16$ and
$\eta_{k-2}(\varepsilon)/6+\bar y\rho^{k-2}\in B_s$. Then $\overline{N\varepsilon+y}\,\rho^k+\nu\notin h+B_s+\Z[i]$ for
every $\nu\in\Lambda_+\setminus\Z[i]$. The same holds with $\rho^{-k}$, $\rho^{2-k}$, $\eta_{2-k}(\varepsilon)$ and
$\Lambda_-$ in place of $\rho^k$, $\rho^{k-2}$, $\eta_{k-2}(\varepsilon)$ and $\Lambda_+$.
\end{lemma}

\begin{proof}
Replacing $(\varepsilon,y,\nu)$ by $(i\varepsilon,iy,-i\nu)$ multiplies $\overline{N\varepsilon+y}\,\rho^k+\nu$ by
$-i$ and replaces $\eta_j$ by $-i\eta_j$; it preserves $h+B_s+\Z[i]$ (as $-ih\equiv h$), $E_q$, $\Lambda_+$ and the
hypotheses. So we may assume $\eta_k(\varepsilon)=1+i$, and then $\eta_{k-2}(\varepsilon)=-1-i$
(Lemma~\ref{lem:values}). Let $u=\bar y\rho^k$, so that $\bar y\rho^{k-2}=u(-7-24i)/25$. Suppose the conclusion
fails. Then $\frac16+\nu_e+u_e\in[-s,s]+\Z$ for $e=1,2$, where we take $\nu_e\in\{\pm\frac15,\pm\frac25\}$. As
$|u_e|<\frac16$ and $s<\frac15$, the values $\nu_e=\frac15$ and $\nu_e=\frac25$ are impossible ($\frac16+\nu_e+u_e$
then lies in $(\frac15,\frac{11}{15})$), and $\nu_e=-\frac25$ forces $u_e\ge\frac1{15}-\delta$. Of the four non-zero
classes $\pm(\frac25,\frac15)$, $\pm(-\frac15,\frac25)$ of $\Lambda_+/\Z[i]$ only $\nu=(-\frac25,-\frac15)$ remains,
so $u_1\ge\frac1{15}-\delta$. On the other hand the hypothesis at $\rho^{k-2}$ says that both coordinates of
$u(-7-24i)/25$ are at least $-\delta$: $7u_1-24u_2\le25\delta$ and $24u_1+7u_2\le25\delta$. Adding $7$ times the
first to $24$ times the second gives $625u_1\le775\delta$, that is $u_1\le\frac{31}{25}\delta$. Hence
$\frac1{15}\le\frac{56}{25}\delta$, that is $\delta\ge\frac5{168}$, a contradiction. The second statement follows by
complex conjugation, as in Lemma~\ref{lem:gapC}: $\bar\varepsilon\in E_q$, $\bar h\equiv h$, and
$\eta_{k-2}(\bar\varepsilon)=\overline{\eta_{2-k}(\varepsilon)}$ by Lemma~\ref{lem:values}(iii).
\end{proof}

\begin{proposition}\label{prop:probe}
Let $17/56<\theta\le1/3$, $k\ge1$ and $N=5^k$. Then
\[
S_N^\theta=\bigcup_{\varepsilon\in E_c}(N\varepsilon+P_k)\ \cup\ \bigcup_{\varepsilon\in E_q}(N\varepsilon+Y_k(\varepsilon)).
\]
In particular every $c\in S_N^\theta$ is $c=N\varepsilon+x$, either with $\varepsilon\in E_c$ and $|x|\le
s\sqrt{10}/3$, or with $\varepsilon\in E_q$ and $|x|\le\sqrt2\,\delta$.
\end{proposition}

For $\theta=1/3$ we have $s\sqrt{10}/3=\sqrt{10}/18$ and $\delta=0$, and this is Proposition~\ref{prop:shape}.

\begin{proof}
The inclusion $\supseteq$ is Lemma~\ref{lem:values}(iv). Note that $\theta>17/56$ means $s<11/56$ and
$\delta<5/168$, and that then $s(1+\sqrt2)<\frac35$. The sets on the right are invariant under $N\Z[i]$ and under
multiplication by $i$ (which maps $E_c$ and $E_q$ to themselves modulo $\Z[i]$), and so is $S_N^\theta$. We use
induction on $k$. For $\lambda\in\Z[i]$ put $\nu_\pm(\lambda)=\frac N5\bar\lambda\rho^{\pm k}$. Since $\frac
N5\rho^k=\frac{(2+i)^{2k}}5=\frac{(2+i)^{2k-1}}{2-i}$, we have $\nu_\pm(\lambda)\in\Lambda_\pm$, and
$\nu_+(\lambda)\in\Z[i]$ if and only if $2+i$ divides $\lambda$; likewise $\nu_-(\lambda)\in\Z[i]$ if and only if
$2-i$ divides $\lambda$. So both are Gaussian integers if and only if $5\mid\lambda$.

\emph{The case $k=1$.} Let $c\in S_5^\theta$. The condition at $\rho^0$ gives $c\in h+B_s+\Z[i]$, so
$c=5h+x+\lambda$ with $x\in B_s$ and $\lambda\in\Z[i]$, which we take modulo $5$; put $z=\bar x$, so $|z|\le s\sqrt2$.
By Lemma~\ref{lem:values}(ii) the conditions at $\rho^{\pm1}$ read $z\rho^{\pm1}+\nu_\pm(\lambda)\in B_s+\Z[i]$, and
as in the proof of Lemma~\ref{lem:gapC} the shift can be taken to be the representative of $\nu_\pm(\lambda)$
modulo $\Z[i]$ with coordinates in $[-\frac25,\frac25]$. If both $\nu_\pm(\lambda)$ are Gaussian integers, then
$5\mid\lambda$ and $x\in P_1$ (as $|z\rho^{\pm1}|\le s\sqrt2<1-s$). If exactly one is, say $\nu_-(\lambda)$, then
$z\rho^{-1}\in B_s$, and Lemma~\ref{lem:gapC} with $k=1$ contradicts the condition at $\rho$; the other case is
symmetric. In the remaining $16$ classes, multiplication by $i$ leaves four orbits, represented by $\lambda=1$, $2$,
$2+2i$ and $1+i$. Write $(z\rho)_{1,2}=\frac15(3z_1-4z_2,\,4z_1+3z_2)$ and
$(z\rho^{-1})_{1,2}=\frac15(3z_1+4z_2,\,-4z_1+3z_2)$.
\begin{itemize}
\item $\lambda=1$: $\nu_\pm\equiv\frac{-2\mp i}5$, so $(z\rho)_1,(z\rho^{-1})_1\ge\frac25-s$. Their sum is
$\frac65z_1\le\frac65s$, so $\frac45-2s\le\frac65s$ and $s\ge\frac14$.
\item $\lambda=2$: $\nu_\pm\equiv\frac{1\mp2i}5$, so $(z\rho)_1,(z\rho^{-1})_1\le s-\frac15$ and
$(z\rho)_2\ge\frac25-s$. As $7(z\rho)_1+25(z\rho^{-1})_1-24(z\rho)_2=0$, we get $0\le32(s-\frac15)+24(s-\frac25)$, so
$s\ge\frac27$.
\item $\lambda=2+2i$: $\nu_+\equiv\frac{-1+2i}5$ and $\nu_-\equiv\frac{-2+i}5$, so $(z\rho)_2\le s-\frac25$,
$(z\rho^{-1})_1\ge\frac25-s$ and $(z\rho^{-1})_2\le s-\frac15$. As $25(z\rho)_2-24(z\rho^{-1})_1+7(z\rho^{-1})_2=0$, we
get $0\le49(s-\frac25)+7(s-\frac15)$, so $s\ge\frac38$.
\item $\lambda=1+i$: then $c=5\varepsilon+y$ with $\varepsilon=\frac{2+2i}3\in E_q$ and $y=x+\frac{1+i}6$, so
$|y|\le\sqrt2(s+\frac16)<0.52$. By Lemma~\ref{lem:values}(iii), $\bar c\rho^j\in h+\eta_j(\varepsilon)/6+\bar
y\rho^j+\Z[i]$, and as $|\bar y\rho^j|<0.52$ and $\frac16+0.52<1-s$, the condition at $\rho^j$ says
$\eta_j(\varepsilon)/6+\bar y\rho^j\in B_s$. So $y\in Y^5_1(\varepsilon)$.
\end{itemize}
The first three orbits are thus empty, as $s<11/56<\frac14$.

\emph{The step from $k-1$ to $k\ge2$.} Let $c\in S_N^\theta$. As $G_{N/5}\subseteq G_N$, $c\in S_{N/5}^\theta$, and by
induction $c=\frac N5\varepsilon'+x$ with $\varepsilon'\in E_c$ and $x\in P_{k-1}$, or with $\varepsilon'\in E_q$ and
$x\in Y^{N/5}_{k-1}(\varepsilon')$.

If $\varepsilon'=h+t\in E_c$, then $\frac N5\varepsilon'=Nh+\frac N5(t-2-2i)$, so $c=Nh+x+\frac N5\lambda$ with
$\lambda\in\Z[i]$. By Lemma~\ref{lem:values}(ii) the condition at $\rho^{\pm k}$ reads $\bar x\rho^{\pm k}+
\nu_\pm(\lambda)\in B_s+\Z[i]$. Here $|x|\le s\sqrt{10}/3<s\sqrt2$ by Lemma~\ref{lem:radii}, and $\bar
x\rho^{\pm(k-2)}\in B_s$ as $x\in P_{k-1}$. By Lemma~\ref{lem:gapC} both $\nu_\pm(\lambda)$ are Gaussian integers,
so $5\mid\lambda$ and $c\in Nh+x+N\Z[i]$; and $\bar x\rho^{\pm k}\in B_s+\Z[i]$ with $|x|<1-s$ gives $x\in P_k$.

If $\varepsilon'=\frac{\alpha'+\beta'i}3+t'\in E_q$, let $\alpha,\beta\in\{1,2\}$ with $\alpha\equiv2\alpha'$ and
$\beta\equiv2\beta'\pmod3$, and $\varepsilon=\frac{\alpha+\beta i}3$. Then $\frac N5\varepsilon'-N\varepsilon\in\frac
N5\Z[i]$, because $\alpha'-5\alpha\equiv\alpha'-2\alpha\equiv0\pmod3$, and likewise for $\beta$; so
$c=N\varepsilon+x+\frac N5\lambda$ with $\lambda\in\Z[i]$. For $|j|\le k-1$ we have $\frac N5\rho^j\in\Z[i]$, so
$\eta^N_j(\varepsilon)=\eta^{N/5}_j(\varepsilon')$ for $|j|\le k-1$, and $x\in Y^N_{k-1}(\varepsilon)$. In
particular $|x|\le\sqrt2\,\delta<\frac16$, and the hypotheses of Lemma~\ref{lem:gapQ} hold at $\rho^{\pm(k-2)}$. The
condition at $\rho^{\pm k}$ says $\overline{N\varepsilon+x}\,\rho^{\pm k}+\nu_\pm(\lambda)\in h+B_s+\Z[i]$, so
by Lemma~\ref{lem:gapQ} both $\nu_\pm(\lambda)$ are Gaussian integers, $5\mid\lambda$ and $c\in N\varepsilon+x+N\Z[i]$.
Finally $\eta_{\pm k}(\varepsilon)/6+\bar x\rho^{\pm k}\in B_s+\Z[i]$ with $|x|<\frac12$ gives $x\in Y_k(\varepsilon)$.
The bounds on $|x|$ are Lemma~\ref{lem:radii}.
\end{proof}

\begin{proof}[Proof of Theorem~\ref{thm:D}]
Suppose that $F^2$ has a homomorphism to $K_{p/q}$ with $p/q<56/17$; we show that $\chi(F^2)\le3$. If $i\in F$ this
cannot happen, as then $\chi(F^2)=\infty$ (proof of Theorem~\ref{thm:B}). If $p/q\le3$, then $\chi(F^2)\le
\chi(K_{p/q})\le3$. Otherwise $\theta=q/p\in(17/56,1/3)$, and we follow the proof of Lemma~\ref{lem:B1} for a finite
$V\subset T$. Theorem~W$^+$ applies, as $p<4q$, and gives a character with $\xi(U)\subseteq[\theta,1-\theta]+\Z$; so
$\varphi(v)\in S_N^\theta$ for $v\in V$, and Proposition~\ref{prop:probe} writes $\varphi(v)=N\varepsilon_v+x_v$ with
$\varepsilon_v\in E$ and $|x_v|\le r_\theta=\max(s\sqrt{10}/3,\sqrt2\,\delta)$ (so $r_{1/3}=r$).
Lemma~\ref{lem:exact} and its proof hold for such points with $r_\theta$ in place of $r$. So, choosing
$N>6r_\theta\max_v(|D_v|+\sum_j|\mu_{v,j}|)$, we obtain $\beta\in L$ with
$\Tr(\beta v)\in E$ for every $v\in V$: the conclusion of Lemma~\ref{lem:B1}. The proof of Lemma~\ref{lem:B2} uses
the $3$-colouring only through this conclusion, and the rest of the proof of Theorem~\ref{thm:B} only through (c) or
(q). So (a) or (b) holds, and $\chi(F^2)\le3$ by Theorem~\ref{thm:B}. Finally, if $\chi(F^2)\ge4$ and
$\chi_c(F^2)<56/17$, the definition of $\chi_c$ as an infimum gives a homomorphism to some $K_{p/q}$ with
$p/q<56/17$, which we have excluded.

For the second statement let $\theta^*=0.3001611$, with $s^*$, $\delta^*$ its values of $s$, $\delta$. It suffices
that for every $k\ge17$ the conclusion of Proposition~\ref{prop:probe} holds for $\theta^*$, so that every $c\in
S_N^{\theta^*}$ is $N\varepsilon+x$ with $\varepsilon\in E$ and $|x|\le r_{\theta^*}<0.22$: then for
$p/q<1/\theta^*$ the argument above applies with $\theta=q/p$, since $S_N^\theta\subseteq S_N^{\theta^*}$, with
$r_{\theta^*}$ in place of $r_\theta$ and $k\ge17$. This is the computer-assisted step. We computed
$S_N^{\theta^*}$ exactly, as a finite union of rational polygons modulo $N$, level by level, with three programs
written separately (two of them by referees): all three find the same number of components at every level $k\le19$,
and only the components of Proposition~\ref{prop:probe} for $k=17,18,19$. For $k\ge18$ the induction step of
Proposition~\ref{prop:probe} goes through with Lemmas~\ref{lem:gapC} and~\ref{lem:gapQ} replaced by the same
statements under the hypotheses $x\in P_{k-1}$, resp.\ $y\in Y^N_{k-1}(\varepsilon)$, for every shift
$\nu\in\Lambda_\pm\setminus\Z[i]$ (the programs do not use the reduction to the shifts of the hand proofs). The
infimum of the $\theta$ for which these statements hold is an exact linear-programming value. It is non-increasing in
$k$: in the coordinates $\bar x\rho^k$, resp.\ $\bar y\rho^k$, and since $\eta_{k-m}=i^m\eta_k$
(Lemma~\ref{lem:values}), the system at level $k+1$ contains the one at level $k$. All three programs find that it is
at most $0.30005568<\theta^*$ for $k=12$, $13$ and $14$ (two of them up to $k=22$).
\end{proof}

\begin{remark}
(1) For a statement valid for every $k$, the bound $17/56$ in Proposition~\ref{prop:probe} is best possible: for
$\theta=17/56$, $S_N^\theta$ has further isolated points for $N=5$, $25$, $125$ and $625$ ($8$, $20$, $12$ and $4$
of them modulo $N\Z[i]$), at which $\min_\gamma\|\operatorname{Re}(\bar c\gamma)\|=17/56$. But the constant $56/17$
comes from the first levels only: by the same computations, for $\theta>333/1106$ the set $S_{5^5}^\theta$ has only
the components of Proposition~\ref{prop:probe}, and the threshold of the induction step is $3303/10981<333/1106$ for
$4\le k\le6$, hence at most that for $k\ge6$; so the levels $k\le5$ and the threshold at $k=6$ already give
$\chi_c(F^2)\ge1106/333\approx3.3213$.

(2) The method cannot reach beyond $10/3$. For every $k$ the point $c_k=Nh+\frac{\rho^k}5-5^{k-1}(2-i)$ lies in
$S_N^{3/10}$, with $\min_\gamma\|\operatorname{Re}(\bar c_k\gamma)\|=\frac3{10}$: for $-k\le j<k$,
$5^{k-1}(2+i)\rho^j\in\Z[i]$ and $\bar c_k\rho^j\equiv h+\rho^{j-k}/5\pmod{\Z[i]}$, at distance $\frac15$ from
$h$; and $\bar c_k\rho^k\equiv h+\frac15-\frac{2+i}5=h-\frac{1+i}5$, because $(2+i)^{2k+1}\equiv2+i\pmod5$. But its distance to
$NE_c$ is at least $\sqrt5\cdot5^{k-1}-\frac15$, and to $NE_q$ at least $\frac76\,5^{k-1}-\frac15$. So for
$\theta\le3/10$ the probe has characters of a third, $5$-adic, kind. On the plane over a number field the places
above $5$, which split in $F(i)$, should exclude them, but the method would have to be extended to include them and
localise at $5$.

(3) The least $\theta$ above which $S_{5^k}^\theta$ has only the components of Proposition~\ref{prop:probe} is
$17/56$ for $k\le4$, $333/1106$ for $k=5$, $3303/10981$ for $k=6$ and about $0.30016093$ for $k=17$. We expect that it
tends to $3/10$, as do the thresholds of the induction step (about $0.300036$ for $20\le k\le22$); that would give
$\chi_c(F^2)\ge10/3$ in Theorem~\ref{thm:D}.
\end{remark}

\section{Computer checks}\label{sec:checks}

The programs are in the folders \nolinkurl{data/number_fields/three_colours/} and
\nolinkurl{data/number_fields/circular/} of \cite{Repo}, with tests in \nolinkurl{tests/test_three_colours.py}.

\emph{Lemma~\ref{lem:B1} as a test.} For a finite $V$, the possible values $z=6\varphi(v)$ form the lattice
$\mathcal M=(\text{column space})\cap\Z^{2m}$, and the types are conditions on $z$ modulo $2$ and $3$; a program
enumerates $\mathcal M/2\mathcal M\times\mathcal M/3\mathcal M$ and decides whether some $\beta$ exists. A second implementation in PARI/GP
\cite{PARI} uses another lattice routine, and a referee wrote a third, which agreed on all $78$ of its trials. The
real quadratic cases of \cite{FC} come out infeasible; $\Q(\sqrt3)$ and $\Q(\sqrt7)$ stay feasible;
$\Q(\sqrt2,\sqrt7)$ is infeasible already with five vectors, $1$, $(7+4\sqrt2i)/9$, $(23+10\sqrt2i)/27$,
$(1+3\sqrt7i)/8$ and $(9+5\sqrt7i)/16$; the sextic field $\Q(2\cos\frac{2\pi}7,\sqrt7)$ is infeasible with $44$
vectors, the obstruction being that of Lemma~\ref{lem:B5} at its two places above $3$ of residue degree $3$;
$\Q(\sqrt3,\sqrt7)$ stays feasible. A referee tested $23$ random quartic and sextic fields: every field predicted to
be $2$- or $3$-colourable stayed feasible, and infeasibility was found for all three quartic fields and four of the
six sextic fields predicted to need four colours (for the other two the random search found no obstruction; the
test is one-sided). The obstructions appear only when $V$ contains vectors that are large at the places above $2$
and $3$ that split in $F(i)$, as the proof of Lemma~\ref{lem:B2} suggests.

\emph{Exact certificates.} The $50$ vectors $G_{25}\{1,u_2,\bar u_2,u_7,\bar u_7\}$ of
Section~\ref{sec:two} (one per pair $\pm u$), and the $70$ vectors $G_{125}V$ for the five vectors above, admit no
character into $[1/3,2/3]$: branch-and-bound trees over the integer values of the relations, with exact Farkas
vectors at the leaves ($1\,579$ and $11\,357$ nodes), accepted by a separate exact checker, which also verifies that
every vector has length $1$ and rejects corrupted copies. With Theorem~W$^+$ this proves
$\chi(\Q(\sqrt2,\sqrt7)^2)\ge4$ without Proposition~\ref{prop:shape} or Sections~\ref{sec:linear}--\ref{sec:three};
a place above $7$ with residue field $\F_7$ gives $\chi\le4$. The certificates of Theorem~\ref{thm:C} are of the
same kind, for the open interval $(2/7,5/7)$: a branch covers all integers strictly inside the range of its
relation (relations must be nonzero), and a leaf is closed by a Farkas vector valid on the open box. They have
$12\,850$ leaves ($d=11$) and $10\,241$ leaves ($d=35$), and two exact checkers that share no code accept them; the
second one also rebuilds $U=G_{25}V$ from $d$. Both reject corrupted copies. The certificate of
Proposition~\ref{prop:fiftynine} ($d=59$, the interval $(15/53,38/53)$) has $127\,181$ nodes and $105\,761$ leaves;
both checkers accept it.

\emph{The lemmas.} Lemma~\ref{lem:B5} was checked exhaustively for $f=1,3,5$, and by the referees with their own
programs (for $f=3$ over all $3^{12}$ $\F_3$-linear maps); the uniqueness step for odd $f\le79$; Lemma~\ref{lem:B4}
on the finite models $\F_{q^2}[\pi]/(\pi^e)$ ($f=1$, $e\le7$, and $f=3$, $e\le2$); Lemma~\ref{lem:B3} over all
$\mu_4$-stable subsets of $\F_9$. For the criterion, a PARI/GP script computes (a) and (b) from a defining
polynomial; for squarefree $|d|<200$ it reproduces the rule for quadratic fields, and it agrees with the known values
$\chi(\Q(\sqrt3,\sqrt{11})^2)=4$ \cite{Fischer1994,MM} and $\chi(\Q(\sqrt2,\sqrt3)^2)=4$ \cite[Theorem~1]{P4C}. For
Section~\ref{sec:circ}, two programs written separately compute $\kappa_1$ for every $(p,f)$ quoted there, and the
bound $|S(b)|\le2\sqrt q$ and the identity of $S(b)$ with a Kloosterman sum were checked numerically for small
$(p,f)$; Proposition~\ref{prop:levels} was checked on the finite models $O_w/\mathfrak m^k$ for $\Q_3$, $\Q_7$,
$\Q_{11}$, $\Q_3(\sqrt3)$, $\Q_7(\sqrt7)$ and $\Q_{27}$.

\emph{Theorem~\ref{thm:D}.} The programs are in \nolinkurl{data/number_fields/circular/probe/}; two referees wrote
their own implementations, which share no code with them or with each other, in its subfolders \nolinkurl{indep/} and
\nolinkurl{indep2/}. All use exact rational arithmetic and compute $S_N^\theta$ level by level as a union of convex
rational polygons: each piece of level $k-1$, translated by the $25$ elements of $\frac N5\{0,\dots,4\}^2$, is cut
along the strips where the coordinates of $\bar c\rho^{\pm k}$ lie in $[\theta,1-\theta]+\Z$; a direct cell-by-cell
computation agrees with this lifting for $k\le3$. Linear-programming values (the thresholds of the induction step,
the least $s$ in the base case, the maximum of $\min_\gamma\|\operatorname{Re}(\bar c\gamma)\|$ over a component) are
floating-point optima certified by exact primal and dual solutions. The implementations agree on everything quoted in
Section~\ref{sec:gap}: the base case (least $s$ equal to $0$, $11/56$ for the eight classes with exactly one integral
$\nu_\pm$, and $1/4$, $2/7$, $3/8$, $1/6$ for the four orbits with the shifts of the proof; with all shifts the third
value is $1/4$), the polygons $P$ and $X(\varepsilon)$, the sets $S_N^\theta$ for small $k$ at several values of
$\theta$ in $(17/56,1/3]$, the extra points at $\theta=17/56$, the thresholds of the induction step, and the
computations at $333/1106$ and at $\theta^*$ (the same number of components at every level $k\le19$). One referee also
tested membership in $S_N^\theta$ exactly at $350\,000$ random points ($N=5,25$), finding none outside the components
of Proposition~\ref{prop:probe}; the $5$-adic points $c_k$ of the Remark were checked exactly for $k\le15$.

\section{Questions}\label{sec:questions}

\begin{enumerate}
\item \emph{Circular local--global.} Is $\chi_c(F^2)\in\{2,3,7/2\}$ whenever $\chi_c(F^2)<4$? More precisely: is
$\chi_c(F^2)=7/2$ exactly when (a) and (b) fail and some place of $F$ above $7$ has residue degree $1$, and
$\chi_c(F^2)\ge4$ when (a) and (b) fail and no such place exists? By Corollary~\ref{cor:localvalues} this is what a
local--global principle for circular colourings would say; Theorem~\ref{thm:B} with Corollary~\ref{cor:B6} is its
part at $2$ and $3$. For $\Q(\sqrt{23})$, where $7$ splits, so that $\chi_c\le7/2$ by Proposition~\ref{prop:C1}, it
predicts $\chi_c=7/2$; for $\Q(\sqrt{59})$, where $7$ is inert and $\chi=4$ \cite[Corollary~2]{FC}, it predicts $\chi_c=4$.
Theorem~\ref{thm:D} is its part below $56/17$ (below $3.3315$ with the computer-assisted step). The method of
Section~\ref{sec:gap} cannot go beyond $10/3$, because of the $5$-adic characters of the probe (Remark in
Section~\ref{sec:gap}); a proof along these lines of the gap up to $7/2$ would have to include them, for instance by
localising at the places above $5$, which split in $F(i)$.
\item \emph{Four colours.} Is $\chi(F^2)$, in general, the least number of colours of a locally constant colouring
of one completion? (For real fields this is \cite[Question~8]{TC}, which also allows Borel colourings at the real
places.) Theorem~\ref{thm:B} answers this at three colours. The first open case is $\Q(\sqrt{47})$, where
$4\le\chi\le5$ (\cite[Corollary~2]{FC} and \cite{Moorhouse}).
\end{enumerate}

\section*{Acknowledgements}
The author thanks the authors of the works cited, whose questions and results this note builds on.

\begin{thebibliography}{99}
\bibitem{BP} M. Benda and M. Perles, \emph{Colorings of metric spaces}, Geombinatorics \textbf{9} (2000),
113--126 (not seen by us).
\bibitem{Davies} J. Davies, \emph{Chromatic number of spacetime}, arXiv:2308.16885 (2023, revised 2024).
\bibitem{DH} H. Davenport and H. Hasse, \emph{Die Nullstellen der Kongruenzzetafunktionen in gewissen zyklischen
F\"allen}, J. Reine Angew. Math. \textbf{172} (1935), 151--182.
\bibitem{dBE} N. G. de Bruijn and P. Erd\H{o}s, \emph{A colour problem for infinite graphs and a problem in the
theory of relations}, Nederl. Akad. Wetensch. Proc. Ser. A \textbf{54} = Indag. Math. \textbf{13} (1951),
369--373.
\bibitem{deGrey} A. D. N. J. de Grey, \emph{The chromatic number of the plane is at least 5}, Geombinatorics
\textbf{28} (2018), 18--31; arXiv:1804.02385.
\bibitem{DeVos} M. DeVos, J. Ebrahimi, M. Ghebleh, L. Goddyn, B. Mohar and R. Naserasr, \emph{Circular coloring the
plane}, SIAM J. Discrete Math. \textbf{21} (2007), 461--465.
\bibitem{Fischer1990} K. G. Fischer, \emph{Additive $K$-colorable extensions of the rational plane}, Discrete
Math. \textbf{82} (1990), 181--195.
\bibitem{Fischer1994} K. G. Fischer, \emph{A planar geometric graph of chromatic number four}, Congr. Numer.
\textbf{104} (1994), 73--79.
\bibitem{Repo} S. Gal\'an, \emph{The Hadwiger--Nelson problem over number fields}, code, data and formal
proofs, 2026. Archived at Zenodo, \href{https://doi.org/10.5281/zenodo.22976635}{doi:10.5281/zenodo.22976635}
(all versions); source at \url{https://github.com/decalion89/chromatic-number-of-the-plane}; the account of this
note is \texttt{notes/three\_colours\_number\_fields.md} and \texttt{notes/circular\_planes.md} there.
\bibitem{P4C} S. Gal\'an, \emph{A short proof that the planes over $\Q(\sqrt3,\sqrt{11})$ and $\Q(\sqrt2,\sqrt3)$ are
$4$-chromatic},
draft, 2026, \texttt{papers/planes-4-chromatic/} in \cite{Repo}.
\bibitem{FC} S. Gal\'an, \emph{Which real quadratic planes need four colours}, draft, 2026,
\texttt{papers/four-colours/} in \cite{Repo}.
\bibitem{PP} S. Gal\'an, \emph{Colouring the $p$-adic plane}, draft, 2026, \texttt{papers/padic-planes/} in
\cite{Repo}.
\bibitem{TC} S. Gal\'an, \emph{Two-colourable planes over number fields}, draft, 2026,
\texttt{papers/two-colour-planes/} in \cite{Repo}.
\bibitem{Winding} S. Gal\'an, \emph{Circular colourings of abelian Cayley graphs below four come from
characters}, draft, 2026, \texttt{papers/winding/} in \cite{Repo}.
\bibitem{Johnson1987} P. D. Johnson Jr., \emph{Two-colorings of real quadratic extensions of $\Q^2$ that forbid
many distances}, Congr. Numer. \textbf{60} (1987), 51--58 (not seen by us).
\bibitem{LN} R. Lidl and H. Niederreiter, \emph{Finite fields}, 2nd ed., Encyclopedia of Mathematics and its
Applications 20, Cambridge University Press, Cambridge, 1997.
\bibitem{Madore} D. A. Madore, \emph{The Hadwiger--Nelson problem over certain fields}, arXiv:1509.07023 (2015).
\bibitem{MM} MildlyMeticulous (GitHub account), \emph{A $2$-adic obstruction to $5$-chromatic unit-distance
graphs, and the exact chromatic number of the plane over $\Q(\sqrt3,\sqrt{11})$}, code, logs and draft, July 2026,
\url{https://github.com/MildlyMeticulous/hn-2adic-obstruction}.
\bibitem{Moorhouse} G. E. Moorhouse, \emph{On the chromatic numbers of planes}, preliminary draft, revised
3 March 2010, \url{https://www.ericmoorhouse.org/pub/chromatic.pdf}.
\bibitem{PARI} The PARI Group, \emph{PARI/GP version 2.15.4}, Univ. Bordeaux, 2023, \url{https://pari.math.u-bordeaux.fr/}.
\bibitem{Weil} A. Weil, \emph{On some exponential sums}, Proc. Nat. Acad. Sci. U.S.A. \textbf{34} (1948), 204--207.
\bibitem{Woodall} D. R. Woodall, \emph{Distances realized by sets covering the plane}, J. Combin. Theory Ser. A
\textbf{14} (1973), 187--200.
\end{thebibliography}

\end{document}
